% Added by DocsPostProcessor:
\setlength{\parskip}{0.125in}
% \\ Added by DocsPostProcessor.

% Replaced by DocsPostProcessor:
% \documentclass[11pt]{report}
\documentclass[a4paper]{report}
% \\ Replaced by DocsPostProcessor.

\def\bl{\mbox{}\newline\mbox{}\newline{}}

% Added by DocsPostProcessor:
\usepackage{amssymb,latexsym,amsmath,amscd,mathtext,ifthen}
\usepackage[unicode]{hyperref}
\usepackage{listings}
\usepackage{literat}
\usepackage{graphicx}
\usepackage[top = 2.0cm]{geometry}
\usepackage[T1,T2A]{fontenc}
\usepackage[utf8]{inputenc}
\usepackage[english,russian]{babel}

\input glyphtounicode
\pdfgentounicode=1

\setlength{\skip\footins}{0.5cm}
\setlength{\footnotesep}{0.5cm}
% \\ Added by DocsPostProcessor.

\usepackage{ifthen}
\newcommand{\hide}[2]{
\ifthenelse{\equal{#1}{inherited}}%
{}%
{}%
}
\newcommand{\entityintro}[3]{%
  \hbox to \hsize{%
    \vbox{%
      \hbox to .2in{}%
    }%
    {\bf #1}%
    \dotfill\pageref{#2}%
  }
  \makebox[\hsize]{%
    \parbox{.4in}{}%
    \parbox[l]{5in}{%
      \vspace{1mm}\it%
      #3%
      \vspace{1mm}%
    }%
  }%
}
\newcommand{\isep}[0]{%
\setlength{\itemsep}{-.4ex}
}
\newcommand{\sld}[0]{%
\setlength{\topsep}{0em}
\setlength{\partopsep}{0em}
\setlength{\parskip}{0em}
\setlength{\parsep}{-1em}
}
\newcommand{\headref}[3]{%
\ifthenelse{#1 = 1}{%
\addcontentsline{toc}{section}{\hspace{\qquad}\protect\numberline{}{#3}}%
}{}%
\ifthenelse{#1 = 2}{%
\addcontentsline{toc}{subsection}{\hspace{\qquad}\protect\numerline{}{#3}}%
}{}%
\ifthenelse{#1 = 3}{%
\addcontentsline{toc}{subsubsection}{\hspace{\qquad}\protect\numerline{}{#3}}%
}{}%
\label{#3}%
\makebox[\textwidth][l]{#2 #3}%
}%
\newcommand{\membername}[1]{{\it #1}\linebreak}
\newcommand{\divideents}[1]{\vskip -1em\indent\rule{2in}{.5mm}}
\newcommand{\refdefined}[1]{
\expandafter\ifx\csname r@#1\endcsname\relax
\relax\else
{$($ in \ref{#1}, page \pageref{#1}$)$}
\fi}
\newcommand{\startsection}[4]{
\gdef\classname{#2}
\subsection{\label{#3}{\bf {\sc #1} #2}}{
\rule[1em]{\hsize}{4pt}\vskip -1em
\vskip .1in 
#4
}%
}
\newcommand{\startsubsubsection}[2]{
\subsubsection{\sc #1}{%
\rule[1em]{\hsize}{2pt}%
#2}
}
\usepackage{color}

% Replaced by DocsPostProcessor:
% \date{DoNotModifyDateHere}
\date{Сентябрь --- октябрь, 2014}
% \\ Replaced by DocsPostProcessor.


% Removed by DocsPostProcessor:
% \pagestyle{myheadings}
% \\ Removed by DocsPostProcessor.

\addtocontents{toc}{\protect\def\protect\packagename{}}
\addtocontents{toc}{\protect\def\protect\classname{}}
\markboth{\protect\packagename -- \protect\classname}{\protect\packagename -- \protect\classname}
\oddsidemargin 0in
\evensidemargin 0in
% \topmargin -.8in
\chardef\bslash=`\\
\textheight 9.4in
\textwidth 6.5in

% Replaced by DocsPostProcessor:
% \title{DoNotModifyTitleHere}
\title{
\includegraphics{images/raic-2013-logo-white-400x150.png}\\
\vspace{0.50in}
\textsc{\Huge CodeHockey 2014}\\
\vspace{0.50in}
\textsc{\LARGE Правила}\\
\textsc{\small Версия 1.0.2}\\
\vspace{2.75in}
\includegraphics{images/mail-ru-and-ssu-logo.png}
}
% \\ Replaced by DocsPostProcessor.

\begin{document}
\maketitle
\sloppy
\raggedright
\tableofcontents

% Added by DocsPostProcessor:
\markboth{qqq}{{\footnotesize CodeHockey 2014}}
\input Tutorial.tex
\raggedright
\sloppy
% \\ Added by DocsPostProcessor.

\gdef\packagename{}
\gdef\classname{}
\newpage
\def\packagename{model}
\chapter{\bf Package model}{
\vskip -.25in
\hbox to \hsize{\it Package Contents\hfil Page}
\rule{\hsize}{.7mm}
\vskip .13in
\hbox{\bf Classes}
\entityintro{ActionType}{l0}{Возможные действия хоккеиста.}
\entityintro{Game}{l1}{Предоставляет доступ к различным игровым константам.}
\entityintro{Hockeyist}{l2}{Класс, определяющий хоккеиста.}
\entityintro{HockeyistState}{l3}{Состояние хоккеиста.}
\entityintro{HockeyistType}{l4}{Тип хоккеиста.}
\entityintro{Move}{l5}{Стратегия игрока может управлять хоккеистом посредством установки свойств объекта данного класса.}
\entityintro{Player}{l6}{Содержит данные о текущем состоянии игрока.}
\entityintro{Puck}{l7}{Класс, определяющий хоккейную шайбу.}
\entityintro{Unit}{l8}{Базовый класс для определения объектов (\textless \textless юнитов\textgreater \textgreater ) на игровом поле.}
\entityintro{World}{l9}{Этот класс описывает игровой мир.}
\vskip .1in
\rule{\hsize}{.7mm}
\vskip .1in
\newpage
\section{Classes}{
\startsection{Class}{ActionType}{l0}{%
{\small Возможные действия хоккеиста.
 
% Removed by DocsPostProcessor:
% \textless p$/$\textgreater
% \\ Removed by DocsPostProcessor.
 
 Хоккеист может совершить действие, если он не сбит с ног ({\tt HockeyistState.KNOCKED\_DOWN}),
 не отдыхает ({\tt HockeyistState.RESTING}) и уже восстановился после своего предыдущего действия
 (значение {\tt hockeyist.remainingCooldownTicks} равно {\tt 0}).
 
% Removed by DocsPostProcessor:
% \textless p$/$\textgreater
% \\ Removed by DocsPostProcessor.
 
 Если хоккеист замахивается клюшкой ({\tt HockeyistState.SWINGING}), то из действий ему доступны только
 {\tt ActionType.STRIKE} и {\tt ActionType.CANCEL\_STRIKE}.}
\vskip .1in 
\startsubsubsection{Declaration}{
\fbox{\vbox{
\hbox{\vbox{\small public final 
class 
ActionType}}
\noindent\hbox{\vbox{{\bf extends} Enum}}
}}}
\startsubsubsection{Fields}{
\begin{itemize}
\item{
public static final ActionType NONE\begin{itemize}\item{\vskip -.9ex Ничего не делать.}\end{itemize}
}
\item{
public static final ActionType TAKE\_PUCK\begin{itemize}\item{\vskip -.9ex Взять шайбу.
 
% Removed by DocsPostProcessor:
% \textless p$/$\textgreater
% \\ Removed by DocsPostProcessor.
 
 Если хоккеист уже контролирует шайбу, либо шайба находится вне зоны досягаемости клюшки хоккеиста (смотрите
 документацию к значениям {\tt game.stickLength} и {\tt game.stickSector}), то действие игнорируется.
 
% Removed by DocsPostProcessor:
% \textless p$/$\textgreater
% \\ Removed by DocsPostProcessor.
 
 В противном случае хоккеист попытается установить контроль над шайбой и, с определённой вероятностью,
 это сделает ((смотрите документацию к {\tt game.pickUpPuckBaseChance} и {\tt game.takePuckAwayBaseChance})).}\end{itemize}
}
\item{
public static final ActionType SWING\begin{itemize}\item{\vskip -.9ex Замахнуться для удара.
 
% Removed by DocsPostProcessor:
% \textless p$/$\textgreater
% \\ Removed by DocsPostProcessor.
 
 Хоккеист замахивается для увеличения силы удара. Чем больше тиков пройдёт с момента начала замаха до удара,
 тем большее воздействие будет на попавшие под удар объекты. Максимальное количество учитываемых тиков ограничено
 значением {\tt game.maxEffectiveSwingTicks}.}\end{itemize}
}
\item{
public static final ActionType STRIKE\begin{itemize}\item{\vskip -.9ex Ударить.
 
% Removed by DocsPostProcessor:
% \textless p$/$\textgreater
% \\ Removed by DocsPostProcessor.
 
 Хоккеист наносит размашистый удар по всем объектам, находящимся в зоне досягаемости его клюшки. Удар может быть
 совершён как с предварительным замахом ({\tt SWING}), так и без него (в этом случае сила удара будет меньше).
 
% Removed by DocsPostProcessor:
% \textless p$/$\textgreater
% \\ Removed by DocsPostProcessor.
 
 Объекты (шайба и хоккеисты, кроме вратарей), попавшие под удар, приобретут некоторый импульс в направлении,
 совпадающим с направлением удара. При ударе по хоккеисту есть также некоторый шанс сбить его с ног.}\end{itemize}
}
\item{
public static final ActionType CANCEL\_STRIKE\begin{itemize}\item{\vskip -.9ex Отменить удар.
 
% Removed by DocsPostProcessor:
% \textless p$/$\textgreater
% \\ Removed by DocsPostProcessor.
 
 Хоккеист выходит из состояния замаха ({\tt SWING}), не совершая удар. Это позволяет
 сэкономить немного выносливости, а также быстрее совершить новое действие
 (смотрите документацию к {\tt game.cancelStrikeActionCooldownTicks}).
 
% Removed by DocsPostProcessor:
% \textless p$/$\textgreater
% \\ Removed by DocsPostProcessor.
 
 Если хоккеист не совершает замах клюшкой, то действие игнорируется.}\end{itemize}
}
\item{
public static final ActionType PASS\begin{itemize}\item{\vskip -.9ex Отдать пас.
 
% Removed by DocsPostProcessor:
% \textless p$/$\textgreater
% \\ Removed by DocsPostProcessor.
 
 Хоккеист пытается передать контролируемую им шайбу другому хоккеисту. Для этого необходимо указать относительную
 силу паса ({\tt move.passPower}) и его направление ({\tt move.passAngle}). В противном случае пас будет отдан
 в направлении, соответствующем направлению хоккеиста, с максимально возможной силой.
 
% Removed by DocsPostProcessor:
% \textless p$/$\textgreater
% \\ Removed by DocsPostProcessor.
 
 Если хоккеист не контролирует шайбу, то действие игнорируется.}\end{itemize}
}
\item{
public static final ActionType SUBSTITUTE\begin{itemize}\item{\vskip -.9ex Заменить активного хоккеиста сидящим на скамейке запасных.
 
% Removed by DocsPostProcessor:
% \textless p$/$\textgreater
% \\ Removed by DocsPostProcessor.
 
 Замена выполняется только на своей половине поля, при этом расстояние от центра хоккеиста до верхней границы
 игровой площадки не должно превышать {\tt game.substitutionAreaHeight}. Дополнительно нужно указать индекс
 хоккеиста ({\tt move.teammateIndex}), на которого будет произведена замена.
 
% Removed by DocsPostProcessor:
% \textless p$/$\textgreater
% \\ Removed by DocsPostProcessor.
 
 Если указан некорректный индекс, или скорость хоккеиста превышает {\tt game.maxSpeedToAllowSubstitute},
 то действие будет проигнорировано.}\end{itemize}
}
\end{itemize}
}
\startsubsubsection{Methods}{
\vskip -2em
\begin{itemize}
\item{\vskip -1.9ex 
\membername{valueOf}
{\tt public static ActionType {\bf valueOf}( {\tt String } {\bf name} )
\label{l10}\label{l11}}%end signature
}%end item
\divideents{values}
\item{\vskip -1.9ex 
\membername{values}
{\tt public static ActionType[] {\bf values}(  )
\label{l12}\label{l13}}%end signature
}%end item
\end{itemize}
}
\hide{inherited}{
\startsubsubsection{Methods inherited from class {\tt Enum}}{
\par{\small 
\refdefined{l14}\vskip -2em
\begin{itemize}
\item{\vskip -1.9ex 
\membername{clone}
{\tt protected final Object {\bf clone}(  )
}%end signature
}%end item
\divideents{compareTo}
\item{\vskip -1.9ex 
\membername{compareTo}
{\tt public final int {\bf compareTo}( {\tt Enum } {\bf arg0} )
}%end signature
}%end item
\divideents{equals}
\item{\vskip -1.9ex 
\membername{equals}
{\tt public final boolean {\bf equals}( {\tt Object } {\bf arg0} )
}%end signature
}%end item
\divideents{finalize}
\item{\vskip -1.9ex 
\membername{finalize}
{\tt protected final void {\bf finalize}(  )
}%end signature
}%end item
\divideents{getDeclaringClass}
\item{\vskip -1.9ex 
\membername{getDeclaringClass}
{\tt public final Class {\bf getDeclaringClass}(  )
}%end signature
}%end item
\divideents{hashCode}
\item{\vskip -1.9ex 
\membername{hashCode}
{\tt public final int {\bf hashCode}(  )
}%end signature
}%end item
\divideents{name}
\item{\vskip -1.9ex 
\membername{name}
{\tt public final String {\bf name}(  )
}%end signature
}%end item
\divideents{ordinal}
\item{\vskip -1.9ex 
\membername{ordinal}
{\tt public final int {\bf ordinal}(  )
}%end signature
}%end item
\divideents{toString}
\item{\vskip -1.9ex 
\membername{toString}
{\tt public String {\bf toString}(  )
}%end signature
}%end item
\divideents{valueOf}
\item{\vskip -1.9ex 
\membername{valueOf}
{\tt public static Enum {\bf valueOf}( {\tt Class } {\bf arg0},
{\tt String } {\bf arg1} )
}%end signature
}%end item
\end{itemize}
}}
}
}
\startsection{Class}{Game}{l1}{%
{\small Предоставляет доступ к различным игровым константам.}
\vskip .1in 
\startsubsubsection{Declaration}{
\fbox{\vbox{
\hbox{\vbox{\small public 
class 
Game}}
\noindent\hbox{\vbox{{\bf extends} Object}}
}}}

% Removed by DocsPostProcessor:
% \startsubsubsection{Constructors}{
% \vskip -2em
% \begin{itemize}
% \item{\vskip -1.9ex 
% \membername{Game}
% {\tt public {\bf Game}( {\tt long } {\bf randomSeed},
% {\tt int } {\bf tickCount},
% {\tt double } {\bf worldWidth},
% {\tt double } {\bf worldHeight},
% {\tt double } {\bf goalNetTop},
% {\tt double } {\bf goalNetWidth},
% {\tt double } {\bf goalNetHeight},
% {\tt double } {\bf rinkTop},
% {\tt double } {\bf rinkLeft},
% {\tt double } {\bf rinkBottom},
% {\tt double } {\bf rinkRight},
% {\tt int } {\bf afterGoalStateTickCount},
% {\tt int } {\bf overtimeTickCount},
% {\tt int } {\bf defaultActionCooldownTicks},
% {\tt int } {\bf swingActionCooldownTicks},
% {\tt int } {\bf cancelStrikeActionCooldownTicks},
% {\tt int } {\bf actionCooldownTicksAfterLosingPuck},
% {\tt double } {\bf stickLength},
% {\tt double } {\bf stickSector},
% {\tt double } {\bf passSector},
% {\tt int } {\bf hockeyistAttributeBaseValue},
% {\tt double } {\bf minActionChance},
% {\tt double } {\bf maxActionChance},
% {\tt double } {\bf strikeAngleDeviation},
% {\tt double } {\bf passAngleDeviation},
% {\tt double } {\bf pickUpPuckBaseChance},
% {\tt double } {\bf takePuckAwayBaseChance},
% {\tt int } {\bf maxEffectiveSwingTicks},
% {\tt double } {\bf strikePowerBaseFactor},
% {\tt double } {\bf strikePowerGrowthFactor},
% {\tt double } {\bf strikePuckBaseChance},
% {\tt double } {\bf knockdownChanceFactor},
% {\tt double } {\bf knockdownTicksFactor},
% {\tt double } {\bf maxSpeedToAllowSubstitute},
% {\tt double } {\bf substitutionAreaHeight},
% {\tt double } {\bf passPowerFactor},
% {\tt double } {\bf hockeyistMaxStamina},
% {\tt double } {\bf activeHockeyistStaminaGrowthPerTick},
% {\tt double } {\bf restingHockeyistStaminaGrowthPerTick},
% {\tt double } {\bf zeroStaminaHockeyistEffectivenessFactor},
% {\tt double } {\bf speedUpStaminaCostFactor},
% {\tt double } {\bf turnStaminaCostFactor},
% {\tt double } {\bf takePuckStaminaCost},
% {\tt double } {\bf swingStaminaCost},
% {\tt double } {\bf strikeStaminaBaseCost},
% {\tt double } {\bf strikeStaminaCostGrowthFactor},
% {\tt double } {\bf cancelStrikeStaminaCost},
% {\tt double } {\bf passStaminaCost},
% {\tt double } {\bf goalieMaxSpeed},
% {\tt double } {\bf hockeyistMaxSpeed},
% {\tt double } {\bf struckHockeyistInitialSpeedFactor},
% {\tt double } {\bf hockeyistSpeedUpFactor},
% {\tt double } {\bf hockeyistSpeedDownFactor},
% {\tt double } {\bf hockeyistTurnAngleFactor},
% {\tt int } {\bf versatileHockeyistStrength},
% {\tt int } {\bf versatileHockeyistEndurance},
% {\tt int } {\bf versatileHockeyistDexterity},
% {\tt int } {\bf versatileHockeyistAgility},
% {\tt int } {\bf forwardHockeyistStrength},
% {\tt int } {\bf forwardHockeyistEndurance},
% {\tt int } {\bf forwardHockeyistDexterity},
% {\tt int } {\bf forwardHockeyistAgility},
% {\tt int } {\bf defencemanHockeyistStrength},
% {\tt int } {\bf defencemanHockeyistEndurance},
% {\tt int } {\bf defencemanHockeyistDexterity},
% {\tt int } {\bf defencemanHockeyistAgility},
% {\tt int } {\bf minRandomHockeyistParameter},
% {\tt int } {\bf maxRandomHockeyistParameter},
% {\tt double } {\bf struckPuckInitialSpeedFactor},
% {\tt double } {\bf puckBindingRange} )
% \label{l15}\label{l16}}%end signature
% }%end item
% \end{itemize}
% }
% \\ Removed by DocsPostProcessor.

\startsubsubsection{Methods}{
\vskip -2em
\begin{itemize}
\item{\vskip -1.9ex 
\membername{getActionCooldownTicksAfterLosingPuck}
{\tt public int {\bf getActionCooldownTicksAfterLosingPuck}(  )
\label{l17}\label{l18}}%end signature
\begin{itemize}
\sld
\item{{\bf Returns} - 
Возвращает длительность задержки, применяемой к хоккеисту
         в случае потери шайбы вследствие воздействия других хоккеистов.
         В течение этого времени хоккеист не может совершать действия. 
}%end item
\end{itemize}
}%end item
\divideents{getActiveHockeyistStaminaGrowthPerTick}
\item{\vskip -1.9ex 
\membername{getActiveHockeyistStaminaGrowthPerTick}
{\tt public double {\bf getActiveHockeyistStaminaGrowthPerTick}(  )
\label{l19}\label{l20}}%end signature
\begin{itemize}
\sld
\item{{\bf Returns} - 
Возвращает значение, на которое увеличивается выносливость хоккеиста за каждый тик в состоянии
         {\tt HockeyistType.ACTIVE}. 
}%end item
\end{itemize}
}%end item
\divideents{getAfterGoalStateTickCount}
\item{\vskip -1.9ex 
\membername{getAfterGoalStateTickCount}
{\tt public int {\bf getAfterGoalStateTickCount}(  )
\label{l21}\label{l22}}%end signature
\begin{itemize}
\sld
\item{{\bf Returns} - 
Возвращает длительность состояния вне игры после гола.
         В течение этого времени новые забитые голы игнорируются,
         а действия не требуют затрат выносливости. 
}%end item
\end{itemize}
}%end item
\divideents{getCancelStrikeActionCooldownTicks}
\item{\vskip -1.9ex 
\membername{getCancelStrikeActionCooldownTicks}
{\tt public int {\bf getCancelStrikeActionCooldownTicks}(  )
\label{l23}\label{l24}}%end signature
\begin{itemize}
\sld
\item{{\bf Returns} - 
Возвращает длительность задержки, применяемой к хоккеисту
         после отмены им удара ({\tt ActionType.CANCEL\_STRIKE}).
         В течение этого времени хоккеист не может совершать новые действия. 
}%end item
\end{itemize}
}%end item
\divideents{getCancelStrikeStaminaCost}
\item{\vskip -1.9ex 
\membername{getCancelStrikeStaminaCost}
{\tt public double {\bf getCancelStrikeStaminaCost}(  )
\label{l25}\label{l26}}%end signature
\begin{itemize}
\sld
\item{{\bf Returns} - 
Возвращает количество выносливости, которое необходимо затратить на совершение действия
         {\tt ActionType.CANCEL\_STRIKE}. 
}%end item
\end{itemize}
}%end item
\divideents{getDefaultActionCooldownTicks}
\item{\vskip -1.9ex 
\membername{getDefaultActionCooldownTicks}
{\tt public int {\bf getDefaultActionCooldownTicks}(  )
\label{l27}\label{l28}}%end signature
\begin{itemize}
\sld
\item{{\bf Returns} - 
Возвращает длительность задержки, применяемой к хоккеисту
         после совершения им большинства действий ({\tt move.action}).
         В течение этого времени хоккеист не может совершать новые действия. 
}%end item
\end{itemize}
}%end item
\divideents{getDefencemanHockeyistAgility}
\item{\vskip -1.9ex 
\membername{getDefencemanHockeyistAgility}
{\tt public int {\bf getDefencemanHockeyistAgility}(  )
\label{l29}\label{l30}}%end signature
\begin{itemize}
\sld
\item{{\bf Returns} - 
Возвращает значение атрибута подвижность для защитника. 
}%end item
\end{itemize}
}%end item
\divideents{getDefencemanHockeyistDexterity}
\item{\vskip -1.9ex 
\membername{getDefencemanHockeyistDexterity}
{\tt public int {\bf getDefencemanHockeyistDexterity}(  )
\label{l31}\label{l32}}%end signature
\begin{itemize}
\sld
\item{{\bf Returns} - 
Возвращает значение атрибута ловкость для защитника. 
}%end item
\end{itemize}
}%end item
\divideents{getDefencemanHockeyistEndurance}
\item{\vskip -1.9ex 
\membername{getDefencemanHockeyistEndurance}
{\tt public int {\bf getDefencemanHockeyistEndurance}(  )
\label{l33}\label{l34}}%end signature
\begin{itemize}
\sld
\item{{\bf Returns} - 
Возвращает значение атрибута стойкость для защитника. 
}%end item
\end{itemize}
}%end item
\divideents{getDefencemanHockeyistStrength}
\item{\vskip -1.9ex 
\membername{getDefencemanHockeyistStrength}
{\tt public int {\bf getDefencemanHockeyistStrength}(  )
\label{l35}\label{l36}}%end signature
\begin{itemize}
\sld
\item{{\bf Returns} - 
Возвращает значение атрибута сила для защитника. 
}%end item
\end{itemize}
}%end item
\divideents{getForwardHockeyistAgility}
\item{\vskip -1.9ex 
\membername{getForwardHockeyistAgility}
{\tt public int {\bf getForwardHockeyistAgility}(  )
\label{l37}\label{l38}}%end signature
\begin{itemize}
\sld
\item{{\bf Returns} - 
Возвращает значение атрибута подвижность для нападающего. 
}%end item
\end{itemize}
}%end item
\divideents{getForwardHockeyistDexterity}
\item{\vskip -1.9ex 
\membername{getForwardHockeyistDexterity}
{\tt public int {\bf getForwardHockeyistDexterity}(  )
\label{l39}\label{l40}}%end signature
\begin{itemize}
\sld
\item{{\bf Returns} - 
Возвращает значение атрибута ловкость для нападающего. 
}%end item
\end{itemize}
}%end item
\divideents{getForwardHockeyistEndurance}
\item{\vskip -1.9ex 
\membername{getForwardHockeyistEndurance}
{\tt public int {\bf getForwardHockeyistEndurance}(  )
\label{l41}\label{l42}}%end signature
\begin{itemize}
\sld
\item{{\bf Returns} - 
Возвращает значение атрибута стойкость для нападающего. 
}%end item
\end{itemize}
}%end item
\divideents{getForwardHockeyistStrength}
\item{\vskip -1.9ex 
\membername{getForwardHockeyistStrength}
{\tt public int {\bf getForwardHockeyistStrength}(  )
\label{l43}\label{l44}}%end signature
\begin{itemize}
\sld
\item{{\bf Returns} - 
Возвращает значение атрибута сила для нападающего. 
}%end item
\end{itemize}
}%end item
\divideents{getGoalieMaxSpeed}
\item{\vskip -1.9ex 
\membername{getGoalieMaxSpeed}
{\tt public double {\bf getGoalieMaxSpeed}(  )
\label{l45}\label{l46}}%end signature
\begin{itemize}
\sld
\item{{\bf Returns} - 
Возвращает максимальную скорость перемещения вратаря. 
}%end item
\end{itemize}
}%end item
\divideents{getGoalNetHeight}
\item{\vskip -1.9ex 
\membername{getGoalNetHeight}
{\tt public double {\bf getGoalNetHeight}(  )
\label{l47}\label{l48}}%end signature
\begin{itemize}
\sld
\item{{\bf Returns} - 
Возвращает высоту ворот. 
}%end item
\end{itemize}
}%end item
\divideents{getGoalNetTop}
\item{\vskip -1.9ex 
\membername{getGoalNetTop}
{\tt public double {\bf getGoalNetTop}(  )
\label{l49}\label{l50}}%end signature
\begin{itemize}
\sld
\item{{\bf Returns} - 
Возвращает ординату верхней штанги ворот. 
}%end item
\end{itemize}
}%end item
\divideents{getGoalNetWidth}
\item{\vskip -1.9ex 
\membername{getGoalNetWidth}
{\tt public double {\bf getGoalNetWidth}(  )
\label{l51}\label{l52}}%end signature
\begin{itemize}
\sld
\item{{\bf Returns} - 
Возвращает ширину ворот. 
}%end item
\end{itemize}
}%end item
\divideents{getHockeyistAttributeBaseValue}
\item{\vskip -1.9ex 
\membername{getHockeyistAttributeBaseValue}
{\tt public int {\bf getHockeyistAttributeBaseValue}(  )
\label{l53}\label{l54}}%end signature
\begin{itemize}
\sld
\item{{\bf Returns} - 
Возвращает базовое значение атрибута хоккеиста.
         Данная величина используется как коэффициент в различных игровых формулах. 
}%end item
\end{itemize}
}%end item
\divideents{getHockeyistMaxSpeed}
\item{\vskip -1.9ex 
\membername{getHockeyistMaxSpeed}
{\tt public double {\bf getHockeyistMaxSpeed}(  )
\label{l55}\label{l56}}%end signature
\begin{itemize}
\sld
\item{{\bf Returns} - 
Возвращает максимальную скорость перемещения полевого хоккеиста. 
}%end item
\end{itemize}
}%end item
\divideents{getHockeyistMaxStamina}
\item{\vskip -1.9ex 
\membername{getHockeyistMaxStamina}
{\tt public double {\bf getHockeyistMaxStamina}(  )
\label{l57}\label{l58}}%end signature
\begin{itemize}
\sld
\item{{\bf Returns} - 
Возвращает максимальное значение выносливости хоккеиста. Выносливость тратится на перемещение
         и совершение хоккеистом различных действий. Каждый тик может восстановиться небольшое количество
         выносливости в зависимости от состояния хоккеиста ({\tt hockeyist.state}). По мере расходования
         выносливости все атрибуты (соответственно, и эффективность всех действий) хоккеиста равномерно
         уменьшаются и достигают значения {\tt zeroStaminaHockeyistEffectivenessFactor} (от начальных
         показателей) при падении выносливости до нуля. Хоккеист не восстанавливает выносливость в состояниях
         {\tt HockeyistState.SWINGING} и {\tt HockeyistState.KNOCKED\_DOWN}. 
}%end item
\end{itemize}
}%end item
\divideents{getHockeyistSpeedDownFactor}
\item{\vskip -1.9ex 
\membername{getHockeyistSpeedDownFactor}
{\tt public double {\bf getHockeyistSpeedDownFactor}(  )
\label{l59}\label{l60}}%end signature
\begin{itemize}
\sld
\item{{\bf Returns} - 
Возвращает модуль ускорения, приобретаемого хоккеистом, при {\tt move.speedUp} равном -1.0,
         базовом значении атрибута подвижность и максимальном запасе выносливости.
         Направление ускорения противоположно направлению хоккеиста.
         В игре отсутствует специальное ограничение на максимальную скорость хоккеиста, однако все
         игровые объекты подвержены воздействию силы трения, которая уменьшает модуль их скорости каждый тик.
         Чем больше скорость, тем на большую величину она уменьшается. 
}%end item
\end{itemize}
}%end item
\divideents{getHockeyistSpeedUpFactor}
\item{\vskip -1.9ex 
\membername{getHockeyistSpeedUpFactor}
{\tt public double {\bf getHockeyistSpeedUpFactor}(  )
\label{l61}\label{l62}}%end signature
\begin{itemize}
\sld
\item{{\bf Returns} - 
Возвращает модуль ускорения, приобретаемого хоккеистом, при {\tt move.speedUp} равном 1.0,
         базовом значении атрибута подвижность и максимальном запасе выносливости.
         Направление ускорения совпадает с направлением хоккеиста.
         В игре отсутствует специальное ограничение на максимальную скорость хоккеиста, однако все
         игровые объекты подвержены воздействию силы трения, которая уменьшает модуль их скорости каждый тик.
         Чем больше скорость, тем на большую величину она уменьшается. 
}%end item
\end{itemize}
}%end item
\divideents{getHockeyistTurnAngleFactor}
\item{\vskip -1.9ex 
\membername{getHockeyistTurnAngleFactor}
{\tt public double {\bf getHockeyistTurnAngleFactor}(  )
\label{l63}\label{l64}}%end signature
\begin{itemize}
\sld
\item{{\bf Returns} - 
Возвращает максимальный модуль угла поворота хоккеиста за тик при базовом значении атрибута подвижность
         и максимальном запасе выносливости. 
}%end item
\end{itemize}
}%end item
\divideents{getKnockdownChanceFactor}
\item{\vskip -1.9ex 
\membername{getKnockdownChanceFactor}
{\tt public double {\bf getKnockdownChanceFactor}(  )
\label{l65}\label{l66}}%end signature
\begin{itemize}
\sld
\item{{\bf Returns} - 
Возвращает шанс ударом ({\tt ActionType.STRIKE}) сбить с ног другого хоккеиста при максимальной
         длительности замаха. Среднее значение атрибутов сила и ловкость хоккеиста, совершающего удар,
         увеличивает шанс сбить с ног. Значение атрибута стойкость атакуемого хоккеиста уменьшает шанс на падение. 
}%end item
\end{itemize}
}%end item
\divideents{getKnockdownTicksFactor}
\item{\vskip -1.9ex 
\membername{getKnockdownTicksFactor}
{\tt public double {\bf getKnockdownTicksFactor}(  )
\label{l67}\label{l68}}%end signature
\begin{itemize}
\sld
\item{{\bf Returns} - 
Возвращает количество тиков, по прошествии которого хоккеист восстановится после падения при базовом
         значении атрибута подвижность. Чем выше подвижность, тем быстрее восстановление. 
}%end item
\end{itemize}
}%end item
\divideents{getMaxActionChance}
\item{\vskip -1.9ex 
\membername{getMaxActionChance}
{\tt public double {\bf getMaxActionChance}(  )
\label{l69}\label{l70}}%end signature
\begin{itemize}
\sld
\item{{\bf Returns} - 
Возвращает максимальный шанс на совершение любого вероятностного действия. 
}%end item
\end{itemize}
}%end item
\divideents{getMaxEffectiveSwingTicks}
\item{\vskip -1.9ex 
\membername{getMaxEffectiveSwingTicks}
{\tt public int {\bf getMaxEffectiveSwingTicks}(  )
\label{l71}\label{l72}}%end signature
\begin{itemize}
\sld
\item{{\bf Returns} - 
Возвращает длительность замаха, после достижения которой сила удара не увеличивается. 
}%end item
\end{itemize}
}%end item
\divideents{getMaxRandomHockeyistParameter}
\item{\vskip -1.9ex 
\membername{getMaxRandomHockeyistParameter}
{\tt public int {\bf getMaxRandomHockeyistParameter}(  )
\label{l73}\label{l74}}%end signature
\begin{itemize}
\sld
\item{{\bf Returns} - 
Возвращает максимально возможное значение любого атрибута для хоккеиста со случайными параметрами. 
}%end item
\end{itemize}
}%end item
\divideents{getMaxSpeedToAllowSubstitute}
\item{\vskip -1.9ex 
\membername{getMaxSpeedToAllowSubstitute}
{\tt public double {\bf getMaxSpeedToAllowSubstitute}(  )
\label{l75}\label{l76}}%end signature
\begin{itemize}
\sld
\item{{\bf Returns} - 
Возвращает максимальную допустимую скорость для выполнения замены хоккеиста. 
}%end item
\end{itemize}
}%end item
\divideents{getMinActionChance}
\item{\vskip -1.9ex 
\membername{getMinActionChance}
{\tt public double {\bf getMinActionChance}(  )
\label{l77}\label{l78}}%end signature
\begin{itemize}
\sld
\item{{\bf Returns} - 
Возвращает минимальный шанс на совершение любого вероятностного действия. 
}%end item
\end{itemize}
}%end item
\divideents{getMinRandomHockeyistParameter}
\item{\vskip -1.9ex 
\membername{getMinRandomHockeyistParameter}
{\tt public int {\bf getMinRandomHockeyistParameter}(  )
\label{l79}\label{l80}}%end signature
\begin{itemize}
\sld
\item{{\bf Returns} - 
Возвращает минимально возможное значение любого атрибута для хоккеиста со случайными параметрами. 
}%end item
\end{itemize}
}%end item
\divideents{getOvertimeTickCount}
\item{\vskip -1.9ex 
\membername{getOvertimeTickCount}
{\tt public int {\bf getOvertimeTickCount}(  )
\label{l81}\label{l82}}%end signature
\begin{itemize}
\sld
\item{{\bf Returns} - 
Возвращает длительность дополнительного времени.
         Дополнительное время наступает в случае ничейного счёта на момент окончания основного времени.
         Если за основное время не было забито ни одного гола, вратари обоих игроков убираются с поля. 
}%end item
\end{itemize}
}%end item
\divideents{getPassAngleDeviation}
\item{\vskip -1.9ex 
\membername{getPassAngleDeviation}
{\tt public double {\bf getPassAngleDeviation}(  )
\label{l83}\label{l84}}%end signature
\begin{itemize}
\sld
\item{{\bf Returns} - 
Возвращает стандартное отклонение распределения Гаусса для угла паса ({\tt ActionType.PASS})
         хоккеиста при базовом значении атрибута ловкость. Чем выше ловкость конкретного хоккеиста,
         тем точнее его пас. 
}%end item
\end{itemize}
}%end item
\divideents{getPassPowerFactor}
\item{\vskip -1.9ex 
\membername{getPassPowerFactor}
{\tt public double {\bf getPassPowerFactor}(  )
\label{l85}\label{l86}}%end signature
\begin{itemize}
\sld
\item{{\bf Returns} - 
Возвращает коэффициент силы паса. Умножается на устанавливаемое стратегией в интервале
         [{\tt 0.0}, {\tt 1.0}] значение силы паса ({\tt move.passPower}). 
}%end item
\end{itemize}
}%end item
\divideents{getPassSector}
\item{\vskip -1.9ex 
\membername{getPassSector}
{\tt public double {\bf getPassSector}(  )
\label{l87}\label{l88}}%end signature
\begin{itemize}
\sld
\item{{\bf Returns} - 
Возвращает сектор, ограничивающий направление паса. 
}%end item
\end{itemize}
}%end item
\divideents{getPassStaminaCost}
\item{\vskip -1.9ex 
\membername{getPassStaminaCost}
{\tt public double {\bf getPassStaminaCost}(  )
\label{l89}\label{l90}}%end signature
\begin{itemize}
\sld
\item{{\bf Returns} - 
Возвращает количество выносливости, которое необходимо затратить на совершение действия
         {\tt ActionType.PASS}. 
}%end item
\end{itemize}
}%end item
\divideents{getPickUpPuckBaseChance}
\item{\vskip -1.9ex 
\membername{getPickUpPuckBaseChance}
{\tt public double {\bf getPickUpPuckBaseChance}(  )
\label{l91}\label{l92}}%end signature
\begin{itemize}
\sld
\item{{\bf Returns} - 
Возвращает шанс подобрать шайбу, не контролируемую другим хоккеистом, без учёта влияния атрибутов
         хоккеиста и скорости шайбы. Равен шансу подобрать шайбу в случае, когда влияющие на действие атрибуты
         хоккеиста равны {\tt hockeyistAttributeBaseValue}, а шайба двигается со скоростью
         {\tt struckPuckInitialSpeedFactor}.
         
% Removed by DocsPostProcessor:
% \textless p$/$\textgreater
% \\ Removed by DocsPostProcessor.
 
         Максимальный из атрибутов ловкость и подвижность хоккеиста увеличивает шанс на захват.
         Скорость шайбы уменьшает шанс на захват. 
}%end item
\end{itemize}
}%end item
\divideents{getPuckBindingRange}
\item{\vskip -1.9ex 
\membername{getPuckBindingRange}
{\tt public double {\bf getPuckBindingRange}(  )
\label{l93}\label{l94}}%end signature
\begin{itemize}
\sld
\item{{\bf Returns} - 
Возвращает расстояние от центра хоккеиста, контролирующего шайбу, до центра шайбы. 
}%end item
\end{itemize}
}%end item
\divideents{getRandomSeed}
\item{\vskip -1.9ex 
\membername{getRandomSeed}
{\tt public long {\bf getRandomSeed}(  )
\label{l95}\label{l96}}%end signature
\begin{itemize}
\sld
\item{{\bf Returns} - 
Возвращает некоторое число, которое ваша стратегия может использовать для инициализации генератора
         случайных чисел. Данное значение имеет рекомендательный характер, однако позволит более точно
         воспроизводить прошедшие игры. 
}%end item
\end{itemize}
}%end item
\divideents{getRestingHockeyistStaminaGrowthPerTick}
\item{\vskip -1.9ex 
\membername{getRestingHockeyistStaminaGrowthPerTick}
{\tt public double {\bf getRestingHockeyistStaminaGrowthPerTick}(  )
\label{l97}\label{l98}}%end signature
\begin{itemize}
\sld
\item{{\bf Returns} - 
Возвращает значение, на которое увеличивается выносливость хоккеиста за каждый тик в состоянии
         {\tt HockeyistType.RESTING}. 
}%end item
\end{itemize}
}%end item
\divideents{getRinkBottom}
\item{\vskip -1.9ex 
\membername{getRinkBottom}
{\tt public double {\bf getRinkBottom}(  )
\label{l99}\label{l100}}%end signature
\begin{itemize}
\sld
\item{{\bf Returns} - 
Возвращает ординату нижней границы игрового поля. 
}%end item
\end{itemize}
}%end item
\divideents{getRinkLeft}
\item{\vskip -1.9ex 
\membername{getRinkLeft}
{\tt public double {\bf getRinkLeft}(  )
\label{l101}\label{l102}}%end signature
\begin{itemize}
\sld
\item{{\bf Returns} - 
Возвращает абсциссу левой границы игрового поля. 
}%end item
\end{itemize}
}%end item
\divideents{getRinkRight}
\item{\vskip -1.9ex 
\membername{getRinkRight}
{\tt public double {\bf getRinkRight}(  )
\label{l103}\label{l104}}%end signature
\begin{itemize}
\sld
\item{{\bf Returns} - 
Возвращает абсциссу правой границы игрового поля. 
}%end item
\end{itemize}
}%end item
\divideents{getRinkTop}
\item{\vskip -1.9ex 
\membername{getRinkTop}
{\tt public double {\bf getRinkTop}(  )
\label{l105}\label{l106}}%end signature
\begin{itemize}
\sld
\item{{\bf Returns} - 
Возвращает ординату верхней границы игрового поля. 
}%end item
\end{itemize}
}%end item
\divideents{getSpeedUpStaminaCostFactor}
\item{\vskip -1.9ex 
\membername{getSpeedUpStaminaCostFactor}
{\tt public double {\bf getSpeedUpStaminaCostFactor}(  )
\label{l107}\label{l108}}%end signature
\begin{itemize}
\sld
\item{{\bf Returns} - 
Возвращает количество выносливости, которое необходимо затратить на максимальное по модулю
         ускорение$/$замедление хоккеиста ({\tt move.speedUp}) за 1 тик. Для меньших значений ускорения затраты
         выносливости пропорционально падают. 
}%end item
\end{itemize}
}%end item
\divideents{getStickLength}
\item{\vskip -1.9ex 
\membername{getStickLength}
{\tt public double {\bf getStickLength}(  )
\label{l109}\label{l110}}%end signature
\begin{itemize}
\sld
\item{{\bf Returns} - 
Возвращает длину клюшки хоккеиста. Хоккеист может воздействовать на игровой объект,
         если и только если расстояние от центра хоккеиста до центра объекта не превышает эту величину. 
}%end item
\end{itemize}
}%end item
\divideents{getStickSector}
\item{\vskip -1.9ex 
\membername{getStickSector}
{\tt public double {\bf getStickSector}(  )
\label{l111}\label{l112}}%end signature
\begin{itemize}
\sld
\item{{\bf Returns} - 
Возвращает сектор клюшки хоккеиста. Хоккеист может воздействовать на игровой объект,
         если и только если угол между направлением хоккеиста и вектором из центра хоккеиста в центр объекта
         не превышает половину этой величины. 
}%end item
\end{itemize}
}%end item
\divideents{getStrikeAngleDeviation}
\item{\vskip -1.9ex 
\membername{getStrikeAngleDeviation}
{\tt public double {\bf getStrikeAngleDeviation}(  )
\label{l113}\label{l114}}%end signature
\begin{itemize}
\sld
\item{{\bf Returns} - 
Возвращает стандартное отклонение распределения Гаусса для угла удара ({\tt ActionType.STRIKE})
         хоккеиста при базовом значении атрибута ловкость. Чем выше ловкость конкретного хоккеиста,
         тем точнее его удар. 
}%end item
\end{itemize}
}%end item
\divideents{getStrikePowerBaseFactor}
\item{\vskip -1.9ex 
\membername{getStrikePowerBaseFactor}
{\tt public double {\bf getStrikePowerBaseFactor}(  )
\label{l115}\label{l116}}%end signature
\begin{itemize}
\sld
\item{{\bf Returns} - 
Возвращает коэффициент силы удара без замаха. 
}%end item
\end{itemize}
}%end item
\divideents{getStrikePowerGrowthFactor}
\item{\vskip -1.9ex 
\membername{getStrikePowerGrowthFactor}
{\tt public double {\bf getStrikePowerGrowthFactor}(  )
\label{l117}\label{l118}}%end signature
\begin{itemize}
\sld
\item{{\bf Returns} - 
Возвращает увеличение коэффициента силы удара за каждый тик замаха.
         Максимальное количество учитываемых тиков ограничено значением {\tt maxEffectiveSwingTicks}. 
}%end item
\end{itemize}
}%end item
\divideents{getStrikePuckBaseChance}
\item{\vskip -1.9ex 
\membername{getStrikePuckBaseChance}
{\tt public double {\bf getStrikePuckBaseChance}(  )
\label{l119}\label{l120}}%end signature
\begin{itemize}
\sld
\item{{\bf Returns} - 
Возвращает базовый шанс ударить шайбу. Базовый шанс не зависит от того,
         контролирует шайбу другой хоккеист или нет, однако на результирующий шанс удара
         по свободной и контролируемой шайбе влияют разные атрибуты хоккеиста
         (смотрите документацию к {\tt pickUpPuckBaseChance} и {\tt takePuckAwayBaseChance}).
         Если хоккеист, совершающий удар, контролирует шайбу, то вероятность удара всегда будет 100\%. 
}%end item
\end{itemize}
}%end item
\divideents{getStrikeStaminaBaseCost}
\item{\vskip -1.9ex 
\membername{getStrikeStaminaBaseCost}
{\tt public double {\bf getStrikeStaminaBaseCost}(  )
\label{l121}\label{l122}}%end signature
\begin{itemize}
\sld
\item{{\bf Returns} - 
Возвращает базовое количество выносливости, которое необходимо затратить на совершение действия
         {\tt ActionType.STRIKE}. 
}%end item
\end{itemize}
}%end item
\divideents{getStrikeStaminaCostGrowthFactor}
\item{\vskip -1.9ex 
\membername{getStrikeStaminaCostGrowthFactor}
{\tt public double {\bf getStrikeStaminaCostGrowthFactor}(  )
\label{l123}\label{l124}}%end signature
\begin{itemize}
\sld
\item{{\bf Returns} - 
Возвращает увеличение затрат выносливости на удар ({\tt ActionType.STRIKE}) за каждый тик замаха.
         Максимальное количество учитываемых тиков ограничено значением {\tt maxEffectiveSwingTicks}. 
}%end item
\end{itemize}
}%end item
\divideents{getStruckHockeyistInitialSpeedFactor}
\item{\vskip -1.9ex 
\membername{getStruckHockeyistInitialSpeedFactor}
{\tt public double {\bf getStruckHockeyistInitialSpeedFactor}(  )
\label{l125}\label{l126}}%end signature
\begin{itemize}
\sld
\item{{\bf Returns} - 
Возвращает модуль скорости, добавляемой хоккеисту, попавшему под удар силы 1.0. 
}%end item
\end{itemize}
}%end item
\divideents{getStruckPuckInitialSpeedFactor}
\item{\vskip -1.9ex 
\membername{getStruckPuckInitialSpeedFactor}
{\tt public double {\bf getStruckPuckInitialSpeedFactor}(  )
\label{l127}\label{l128}}%end signature
\begin{itemize}
\sld
\item{{\bf Returns} - 
Возвращает модуль скорости, устанавливаемой шайбе, попавшей под удар силы 1.0. 
}%end item
\end{itemize}
}%end item
\divideents{getSubstitutionAreaHeight}
\item{\vskip -1.9ex 
\membername{getSubstitutionAreaHeight}
{\tt public double {\bf getSubstitutionAreaHeight}(  )
\label{l129}\label{l130}}%end signature
\begin{itemize}
\sld
\item{{\bf Returns} - 
Возвращает высоту зоны, в которой может быть выполнена замена хоккеиста. Зона расположена вдоль верхней
         границы игровой площадки. Замена может быть выполнена только на своей половине поля. 
}%end item
\end{itemize}
}%end item
\divideents{getSwingActionCooldownTicks}
\item{\vskip -1.9ex 
\membername{getSwingActionCooldownTicks}
{\tt public int {\bf getSwingActionCooldownTicks}(  )
\label{l131}\label{l132}}%end signature
\begin{itemize}
\sld
\item{{\bf Returns} - 
Возвращает длительность задержки, применяемой к хоккеисту
         после совершения им действия замах ({\tt ActionType.SWING}).
         В течение этого времени хоккеист не может совершать новые действия. 
}%end item
\end{itemize}
}%end item
\divideents{getSwingStaminaCost}
\item{\vskip -1.9ex 
\membername{getSwingStaminaCost}
{\tt public double {\bf getSwingStaminaCost}(  )
\label{l133}\label{l134}}%end signature
\begin{itemize}
\sld
\item{{\bf Returns} - 
Возвращает количество выносливости, которое необходимо затратить на совершение действия
         {\tt ActionType.SWING}. 
}%end item
\end{itemize}
}%end item
\divideents{getTakePuckAwayBaseChance}
\item{\vskip -1.9ex 
\membername{getTakePuckAwayBaseChance}
{\tt public double {\bf getTakePuckAwayBaseChance}(  )
\label{l135}\label{l136}}%end signature
\begin{itemize}
\sld
\item{{\bf Returns} - 
Возвращает базовый шанс отнять шайбу у другого хоккеиста.
         
% Removed by DocsPostProcessor:
% \textless p$/$\textgreater
% \\ Removed by DocsPostProcessor.
 
         Максимальный из атрибутов сила и ловкость хоккеиста, отнимающего шайбу, увеличивает шанс на захват.
         Максимальный из атрибутов стойкость и подвижность текущего владельца шайбы уменьшает шанс на её потерю. 
}%end item
\end{itemize}
}%end item
\divideents{getTakePuckStaminaCost}
\item{\vskip -1.9ex 
\membername{getTakePuckStaminaCost}
{\tt public double {\bf getTakePuckStaminaCost}(  )
\label{l137}\label{l138}}%end signature
\begin{itemize}
\sld
\item{{\bf Returns} - 
Возвращает количество выносливости, которое необходимо затратить на совершение действия
         {\tt ActionType.TAKE\_PUCK}. 
}%end item
\end{itemize}
}%end item
\divideents{getTickCount}
\item{\vskip -1.9ex 
\membername{getTickCount}
{\tt public int {\bf getTickCount}(  )
\label{l139}\label{l140}}%end signature
\begin{itemize}
\sld
\item{{\bf Returns} - 
Возвращает длительность игры в тиках. 
}%end item
\end{itemize}
}%end item
\divideents{getTurnStaminaCostFactor}
\item{\vskip -1.9ex 
\membername{getTurnStaminaCostFactor}
{\tt public double {\bf getTurnStaminaCostFactor}(  )
\label{l141}\label{l142}}%end signature
\begin{itemize}
\sld
\item{{\bf Returns} - 
Возвращает количество выносливости, которое необходимо затратить на максимальный по модулю
         угол поворота хоккеиста ({\tt move.turn}) за 1 тик. Для меньших значений угла поворота затраты
         выносливости пропорционально падают. 
}%end item
\end{itemize}
}%end item
\divideents{getVersatileHockeyistAgility}
\item{\vskip -1.9ex 
\membername{getVersatileHockeyistAgility}
{\tt public int {\bf getVersatileHockeyistAgility}(  )
\label{l143}\label{l144}}%end signature
\begin{itemize}
\sld
\item{{\bf Returns} - 
Возвращает значение атрибута подвижность для хоккеиста-универсала. 
}%end item
\end{itemize}
}%end item
\divideents{getVersatileHockeyistDexterity}
\item{\vskip -1.9ex 
\membername{getVersatileHockeyistDexterity}
{\tt public int {\bf getVersatileHockeyistDexterity}(  )
\label{l145}\label{l146}}%end signature
\begin{itemize}
\sld
\item{{\bf Returns} - 
Возвращает значение атрибута ловкость для хоккеиста-универсала. 
}%end item
\end{itemize}
}%end item
\divideents{getVersatileHockeyistEndurance}
\item{\vskip -1.9ex 
\membername{getVersatileHockeyistEndurance}
{\tt public int {\bf getVersatileHockeyistEndurance}(  )
\label{l147}\label{l148}}%end signature
\begin{itemize}
\sld
\item{{\bf Returns} - 
Возвращает значение атрибута стойкость для хоккеиста-универсала. 
}%end item
\end{itemize}
}%end item
\divideents{getVersatileHockeyistStrength}
\item{\vskip -1.9ex 
\membername{getVersatileHockeyistStrength}
{\tt public int {\bf getVersatileHockeyistStrength}(  )
\label{l149}\label{l150}}%end signature
\begin{itemize}
\sld
\item{{\bf Returns} - 
Возвращает значение атрибута сила для хоккеиста-универсала. 
}%end item
\end{itemize}
}%end item
\divideents{getWorldHeight}
\item{\vskip -1.9ex 
\membername{getWorldHeight}
{\tt public double {\bf getWorldHeight}(  )
\label{l151}\label{l152}}%end signature
\begin{itemize}
\sld
\item{{\bf Returns} - 
Возвращает высоту игрового мира. 
}%end item
\end{itemize}
}%end item
\divideents{getWorldWidth}
\item{\vskip -1.9ex 
\membername{getWorldWidth}
{\tt public double {\bf getWorldWidth}(  )
\label{l153}\label{l154}}%end signature
\begin{itemize}
\sld
\item{{\bf Returns} - 
Возвращает ширину игрового мира. 
}%end item
\end{itemize}
}%end item
\divideents{getZeroStaminaHockeyistEffectivenessFactor}
\item{\vskip -1.9ex 
\membername{getZeroStaminaHockeyistEffectivenessFactor}
{\tt public double {\bf getZeroStaminaHockeyistEffectivenessFactor}(  )
\label{l155}\label{l156}}%end signature
\begin{itemize}
\sld
\item{{\bf Returns} - 
Возвращает коэффициент эффективности действий хоккеиста при падении его выносливости до нуля. 
}%end item
\end{itemize}
}%end item
\end{itemize}
}
\hide{inherited}{
}
}
\startsection{Class}{Hockeyist}{l2}{%
{\small Класс, определяющий хоккеиста. Содержит также все свойства юнита.}
\vskip .1in 
\startsubsubsection{Declaration}{
\fbox{\vbox{
\hbox{\vbox{\small public 
class 
Hockeyist}}
\noindent\hbox{\vbox{{\bf extends} Unit}}
}}}

% Removed by DocsPostProcessor:
% \startsubsubsection{Constructors}{
% \vskip -2em
% \begin{itemize}
% \item{\vskip -1.9ex 
% \membername{Hockeyist}
% {\tt public {\bf Hockeyist}( {\tt long } {\bf id},
% {\tt long } {\bf playerId},
% {\tt int } {\bf teammateIndex},
% {\tt double } {\bf mass},
% {\tt double } {\bf radius},
% {\tt double } {\bf x},
% {\tt double } {\bf y},
% {\tt double } {\bf speedX},
% {\tt double } {\bf speedY},
% {\tt double } {\bf angle},
% {\tt double } {\bf angularSpeed},
% {\tt boolean } {\bf teammate},
% {\tt HockeyistType } {\bf type},
% {\tt int } {\bf strength},
% {\tt int } {\bf endurance},
% {\tt int } {\bf dexterity},
% {\tt int } {\bf agility},
% {\tt double } {\bf stamina},
% {\tt HockeyistState } {\bf state},
% {\tt int } {\bf originalPositionIndex},
% {\tt int } {\bf remainingKnockdownTicks},
% {\tt int } {\bf remainingCooldownTicks},
% {\tt int } {\bf swingTicks},
% {\tt ActionType } {\bf lastAction},
% {\tt Integer } {\bf lastActionTick} )
% \label{l157}\label{l158}}%end signature
% }%end item
% \end{itemize}
% }
% \\ Removed by DocsPostProcessor.

\startsubsubsection{Methods}{
\vskip -2em
\begin{itemize}
\item{\vskip -1.9ex 
\membername{getAgility}
{\tt public int {\bf getAgility}(  )
\label{l159}\label{l160}}%end signature
\begin{itemize}
\sld
\item{{\bf Returns} - 
Возвращает значение атрибута подвижность. 
}%end item
\end{itemize}
}%end item
\divideents{getDexterity}
\item{\vskip -1.9ex 
\membername{getDexterity}
{\tt public int {\bf getDexterity}(  )
\label{l161}\label{l162}}%end signature
\begin{itemize}
\sld
\item{{\bf Returns} - 
Возвращает значение атрибута ловкость. 
}%end item
\end{itemize}
}%end item
\divideents{getEndurance}
\item{\vskip -1.9ex 
\membername{getEndurance}
{\tt public int {\bf getEndurance}(  )
\label{l163}\label{l164}}%end signature
\begin{itemize}
\sld
\item{{\bf Returns} - 
Возвращает значение атрибута стойкость. 
}%end item
\end{itemize}
}%end item
\divideents{getLastAction}
\item{\vskip -1.9ex 
\membername{getLastAction}
{\tt public ActionType {\bf getLastAction}(  )
\label{l165}\label{l166}}%end signature
\begin{itemize}
\sld
\item{{\bf Returns} - 
Возвращает последнее действие ({\tt move.action}), совершённое хоккеистом, или {\tt null}
         ({\tt UNKNOWN\_ACTION} в пакетах некоторых языков) в случае, если хоккеист ещё не совершил
         ни одного действия. 
}%end item
\end{itemize}
}%end item
\divideents{getLastActionTick}
\item{\vskip -1.9ex 
\membername{getLastActionTick}
{\tt public Integer {\bf getLastActionTick}(  )
\label{l167}\label{l168}}%end signature
\begin{itemize}
\sld
\item{{\bf Returns} - 
Возвращает номер тика, в который хоккеист совершил своё последние действие ({\tt move.action}),
         или {\tt null} ({\tt -1} в пакетах некоторых языков) в случае, если хоккеист ещё не совершил
         ни одного действия. 
}%end item
\end{itemize}
}%end item
\divideents{getOriginalPositionIndex}
\item{\vskip -1.9ex 
\membername{getOriginalPositionIndex}
{\tt public int {\bf getOriginalPositionIndex}(  )
\label{l169}\label{l170}}%end signature
\begin{itemize}
\sld
\item{{\bf Returns} - 
Возвращает индекс исходной позиции хоккеиста или {\tt -1} для вратаря или хоккеиста,
         отдыхающего за пределами игрового поля. На эту позицию хоккеист будет помещён при разыгрывании шайбы.
         При выполнении действия замена {\tt ActionType.SUBSTITUTE} индексы исходных позиций хоккеистов,
         участвующих в замене, меняются местами. 
}%end item
\end{itemize}
}%end item
\divideents{getPlayerId}
\item{\vskip -1.9ex 
\membername{getPlayerId}
{\tt public long {\bf getPlayerId}(  )
\label{l171}\label{l172}}%end signature
\begin{itemize}
\sld
\item{{\bf Returns} - 
Возвращает идентификатор игрока, в команду которого входит хоккеист. 
}%end item
\end{itemize}
}%end item
\divideents{getRemainingCooldownTicks}
\item{\vskip -1.9ex 
\membername{getRemainingCooldownTicks}
{\tt public int {\bf getRemainingCooldownTicks}(  )
\label{l173}\label{l174}}%end signature
\begin{itemize}
\sld
\item{{\bf Returns} - 
Возвращает количество тиков, по прошествии которого хоккеист сможет совершить какое-либо
         действие ({\tt move.action}), или {\tt 0}, если хоккеист может совершить действие в данный тик. 
}%end item
\end{itemize}
}%end item
\divideents{getRemainingKnockdownTicks}
\item{\vskip -1.9ex 
\membername{getRemainingKnockdownTicks}
{\tt public int {\bf getRemainingKnockdownTicks}(  )
\label{l175}\label{l176}}%end signature
\begin{itemize}
\sld
\item{{\bf Returns} - 
Возвращает количество тиков, по прошествии которого хоккеист восстановится после падения,
         или {\tt 0}, если хоккеист не сбит с ног. 
}%end item
\end{itemize}
}%end item
\divideents{getStamina}
\item{\vskip -1.9ex 
\membername{getStamina}
{\tt public double {\bf getStamina}(  )
\label{l177}\label{l178}}%end signature
\begin{itemize}
\sld
\item{{\bf Returns} - 
Возвращает текущее значение выносливости. 
}%end item
\end{itemize}
}%end item
\divideents{getState}
\item{\vskip -1.9ex 
\membername{getState}
{\tt public HockeyistState {\bf getState}(  )
\label{l179}\label{l180}}%end signature
\begin{itemize}
\sld
\item{{\bf Returns} - 
Возвращает состояние хоккеиста. 
}%end item
\end{itemize}
}%end item
\divideents{getStrength}
\item{\vskip -1.9ex 
\membername{getStrength}
{\tt public int {\bf getStrength}(  )
\label{l181}\label{l182}}%end signature
\begin{itemize}
\sld
\item{{\bf Returns} - 
Возвращает значение атрибута сила. 
}%end item
\end{itemize}
}%end item
\divideents{getSwingTicks}
\item{\vskip -1.9ex 
\membername{getSwingTicks}
{\tt public int {\bf getSwingTicks}(  )
\label{l183}\label{l184}}%end signature
\begin{itemize}
\sld
\item{{\bf Returns} - 
Для хоккеиста, находящегося в состоянии замаха ({\tt HockeyistState.SWINGING}),
         возвращает количество тиков, прошедших от начала замаха. В противном случае возвращает {\tt 0}. 
}%end item
\end{itemize}
}%end item
\divideents{getTeammateIndex}
\item{\vskip -1.9ex 
\membername{getTeammateIndex}
{\tt public int {\bf getTeammateIndex}(  )
\label{l185}\label{l186}}%end signature
\begin{itemize}
\sld
\item{{\bf Returns} - 
Возвращает 0-индексированный номер хоккеиста в команде. 
}%end item
\end{itemize}
}%end item
\divideents{getType}
\item{\vskip -1.9ex 
\membername{getType}
{\tt public HockeyistType {\bf getType}(  )
\label{l187}\label{l188}}%end signature
\begin{itemize}
\sld
\item{{\bf Returns} - 
Возвращает тип хоккеиста. 
}%end item
\end{itemize}
}%end item
\divideents{isTeammate}
\item{\vskip -1.9ex 
\membername{isTeammate}
{\tt public boolean {\bf isTeammate}(  )
\label{l189}\label{l190}}%end signature
\begin{itemize}
\sld
\item{{\bf Returns} - 
Возвращает {\tt true}, если и только если данный хоккеист входит в команду вашего игрока. 
}%end item
\end{itemize}
}%end item
\end{itemize}
}
\hide{inherited}{
\startsubsubsection{Methods inherited from class {\tt Unit}}{
\par{\small 
\refdefined{l8}\vskip -2em
\begin{itemize}
\item{\vskip -1.9ex 
\membername{getAngle}
{\tt public final double {\bf getAngle}(  )
}%end signature
\begin{itemize}
\sld
\item{{\bf Returns} - 
Возвращает угол поворота объекта в радианах. Нулевой угол соответствует направлению оси абсцисс.
         Положительные значения соответствуют повороту по часовой стрелке. 
}%end item
\end{itemize}
}%end item
\divideents{getAngleTo}
\item{\vskip -1.9ex 
\membername{getAngleTo}
{\tt public double {\bf getAngleTo}( {\tt double } {\bf x},
{\tt double } {\bf y} )
}%end signature
\begin{itemize}
\sld
\item{
\sld
{\bf Parameters}
\sld\isep
  \begin{itemize}
\sld\isep
   \item{
\sld
{\tt x} - X-координата точки.}
   \item{
\sld
{\tt y} - Y-координата точки.}
  \end{itemize}
}%end item
\item{{\bf Returns} - 
Возвращает ориентированный угол [{\tt -PI}, {\tt PI}] между направлением
         данного объекта и вектором из центра данного объекта к указанной точке. 
}%end item
\end{itemize}
}%end item
\divideents{getAngleTo}
\item{\vskip -1.9ex 
\membername{getAngleTo}
{\tt public double {\bf getAngleTo}( {\tt Unit } {\bf unit} )
}%end signature
\begin{itemize}
\sld
\item{
\sld
{\bf Parameters}
\sld\isep
  \begin{itemize}
\sld\isep
   \item{
\sld
{\tt unit} - Объект, к центру которого необходимо определить угол.}
  \end{itemize}
}%end item
\item{{\bf Returns} - 
Возвращает ориентированный угол [{\tt -PI}, {\tt PI}] между направлением
         данного объекта и вектором из центра данного объекта к центру указанного объекта. 
}%end item
\end{itemize}
}%end item
\divideents{getAngularSpeed}
\item{\vskip -1.9ex 
\membername{getAngularSpeed}
{\tt public double {\bf getAngularSpeed}(  )
}%end signature
\begin{itemize}
\sld
\item{{\bf Returns} - 
Возвращает скорость вращения объекта.
         Положительные значения соответствуют вращению по часовой стрелке. 
}%end item
\end{itemize}
}%end item
\divideents{getDistanceTo}
\item{\vskip -1.9ex 
\membername{getDistanceTo}
{\tt public double {\bf getDistanceTo}( {\tt double } {\bf x},
{\tt double } {\bf y} )
}%end signature
\begin{itemize}
\sld
\item{
\sld
{\bf Parameters}
\sld\isep
  \begin{itemize}
\sld\isep
   \item{
\sld
{\tt x} - X-координата точки.}
   \item{
\sld
{\tt y} - Y-координата точки.}
  \end{itemize}
}%end item
\item{{\bf Returns} - 
Возвращает расстояние до точки от центра данного объекта. 
}%end item
\end{itemize}
}%end item
\divideents{getDistanceTo}
\item{\vskip -1.9ex 
\membername{getDistanceTo}
{\tt public double {\bf getDistanceTo}( {\tt Unit } {\bf unit} )
}%end signature
\begin{itemize}
\sld
\item{
\sld
{\bf Parameters}
\sld\isep
  \begin{itemize}
\sld\isep
   \item{
\sld
{\tt unit} - Объект, до центра которого необходимо определить расстояние.}
  \end{itemize}
}%end item
\item{{\bf Returns} - 
Возвращает расстояние от центра данного объекта до центра указанного объекта. 
}%end item
\end{itemize}
}%end item
\divideents{getId}
\item{\vskip -1.9ex 
\membername{getId}
{\tt public long {\bf getId}(  )
}%end signature
\begin{itemize}
\sld
\item{{\bf Returns} - 
Возвращает уникальный идентификатор объекта. 
}%end item
\end{itemize}
}%end item
\divideents{getMass}
\item{\vskip -1.9ex 
\membername{getMass}
{\tt public double {\bf getMass}(  )
}%end signature
\begin{itemize}
\sld
\item{{\bf Returns} - 
Возвращает массу объекта в единицах массы. 
}%end item
\end{itemize}
}%end item
\divideents{getRadius}
\item{\vskip -1.9ex 
\membername{getRadius}
{\tt public double {\bf getRadius}(  )
}%end signature
\begin{itemize}
\sld
\item{{\bf Returns} - 
Возвращает радиус объекта. 
}%end item
\end{itemize}
}%end item
\divideents{getSpeedX}
\item{\vskip -1.9ex 
\membername{getSpeedX}
{\tt public final double {\bf getSpeedX}(  )
}%end signature
\begin{itemize}
\sld
\item{{\bf Returns} - 
Возвращает X-составляющую скорости объекта. Ось абсцисс направлена слева направо. 
}%end item
\end{itemize}
}%end item
\divideents{getSpeedY}
\item{\vskip -1.9ex 
\membername{getSpeedY}
{\tt public final double {\bf getSpeedY}(  )
}%end signature
\begin{itemize}
\sld
\item{{\bf Returns} - 
Возвращает Y-составляющую скорости объекта. Ось ординат направлена свеху вниз. 
}%end item
\end{itemize}
}%end item
\divideents{getX}
\item{\vskip -1.9ex 
\membername{getX}
{\tt public final double {\bf getX}(  )
}%end signature
\begin{itemize}
\sld
\item{{\bf Returns} - 
Возвращает X-координату центра объекта. Ось абсцисс направлена слева направо. 
}%end item
\end{itemize}
}%end item
\divideents{getY}
\item{\vskip -1.9ex 
\membername{getY}
{\tt public final double {\bf getY}(  )
}%end signature
\begin{itemize}
\sld
\item{{\bf Returns} - 
Возвращает Y-координату центра объекта. Ось ординат направлена свеху вниз. 
}%end item
\end{itemize}
}%end item
\end{itemize}
}}
}
}
\startsection{Class}{HockeyistState}{l3}{%
{\small Состояние хоккеиста.}
\vskip .1in 
\startsubsubsection{Declaration}{
\fbox{\vbox{
\hbox{\vbox{\small public final 
class 
HockeyistState}}
\noindent\hbox{\vbox{{\bf extends} Enum}}
}}}
\startsubsubsection{Fields}{
\begin{itemize}
\item{
public static final HockeyistState ACTIVE\begin{itemize}\item{\vskip -.9ex Хоккеист находится на игровом поле.}\end{itemize}
}
\item{
public static final HockeyistState SWINGING\begin{itemize}\item{\vskip -.9ex Хоккеист находится на игровом поле и делает замах клюшкой.
 
% Removed by DocsPostProcessor:
% \textless p$/$\textgreater
% \\ Removed by DocsPostProcessor.
 
 Во время замаха стратегия не может управлять движением хоккеиста, а из действий доступны только
 {\tt ActionType.STRIKE} и {\tt ActionType.CANCEL\_STRIKE}.}\end{itemize}
}
\item{
public static final HockeyistState KNOCKED\_DOWN\begin{itemize}\item{\vskip -.9ex Хоккеист находится на игровом поле, но сбит с ног.
 Стратегия игрока не может им управлять.}\end{itemize}
}
\item{
public static final HockeyistState RESTING\begin{itemize}\item{\vskip -.9ex Хоккеист отдыхает вне игрового поля.
 Стратегия игрока не может им управлять.}\end{itemize}
}
\end{itemize}
}
\startsubsubsection{Methods}{
\vskip -2em
\begin{itemize}
\item{\vskip -1.9ex 
\membername{valueOf}
{\tt public static HockeyistState {\bf valueOf}( {\tt String } {\bf name} )
\label{l191}\label{l192}}%end signature
}%end item
\divideents{values}
\item{\vskip -1.9ex 
\membername{values}
{\tt public static HockeyistState[] {\bf values}(  )
\label{l193}\label{l194}}%end signature
}%end item
\end{itemize}
}
\hide{inherited}{
\startsubsubsection{Methods inherited from class {\tt Enum}}{
\par{\small 
\refdefined{l14}\vskip -2em
\begin{itemize}
\item{\vskip -1.9ex 
\membername{clone}
{\tt protected final Object {\bf clone}(  )
}%end signature
}%end item
\divideents{compareTo}
\item{\vskip -1.9ex 
\membername{compareTo}
{\tt public final int {\bf compareTo}( {\tt Enum } {\bf arg0} )
}%end signature
}%end item
\divideents{equals}
\item{\vskip -1.9ex 
\membername{equals}
{\tt public final boolean {\bf equals}( {\tt Object } {\bf arg0} )
}%end signature
}%end item
\divideents{finalize}
\item{\vskip -1.9ex 
\membername{finalize}
{\tt protected final void {\bf finalize}(  )
}%end signature
}%end item
\divideents{getDeclaringClass}
\item{\vskip -1.9ex 
\membername{getDeclaringClass}
{\tt public final Class {\bf getDeclaringClass}(  )
}%end signature
}%end item
\divideents{hashCode}
\item{\vskip -1.9ex 
\membername{hashCode}
{\tt public final int {\bf hashCode}(  )
}%end signature
}%end item
\divideents{name}
\item{\vskip -1.9ex 
\membername{name}
{\tt public final String {\bf name}(  )
}%end signature
}%end item
\divideents{ordinal}
\item{\vskip -1.9ex 
\membername{ordinal}
{\tt public final int {\bf ordinal}(  )
}%end signature
}%end item
\divideents{toString}
\item{\vskip -1.9ex 
\membername{toString}
{\tt public String {\bf toString}(  )
}%end signature
}%end item
\divideents{valueOf}
\item{\vskip -1.9ex 
\membername{valueOf}
{\tt public static Enum {\bf valueOf}( {\tt Class } {\bf arg0},
{\tt String } {\bf arg1} )
}%end signature
}%end item
\end{itemize}
}}
}
}
\startsection{Class}{HockeyistType}{l4}{%
{\small Тип хоккеиста.
 
% Removed by DocsPostProcessor:
% \textless p$/$\textgreater
% \\ Removed by DocsPostProcessor.
 
 В Раунде 1 чемпионата стратегия игрока управляет двумя хоккеистами-универсалами ({\tt VERSATILE}).
 
% Removed by DocsPostProcessor:
% \textless p$/$\textgreater
% \\ Removed by DocsPostProcessor.
 
 В Раунде 2 команда состоит из одного {\tt VERSATILE}, одного {\tt FORWARD} и одного {\tt DEFENCEMAN}.
 
% Removed by DocsPostProcessor:
% \textless p$/$\textgreater
% \\ Removed by DocsPostProcessor.
 
 В Групповом отборе и Финале у игрока в распоряжении имеется шесть хоккеистов типа {\tt RANDOM},
 однако одновременно в игре могут участвовать только трое из них. Остальные должны находиться на скамейке запасных.}
\vskip .1in 
\startsubsubsection{Declaration}{
\fbox{\vbox{
\hbox{\vbox{\small public final 
class 
HockeyistType}}
\noindent\hbox{\vbox{{\bf extends} Enum}}
}}}
\startsubsubsection{Fields}{
\begin{itemize}
\item{
public static final HockeyistType GOALIE\begin{itemize}\item{\vskip -.9ex Вратарь.
 
% Removed by DocsPostProcessor:
% \textless p$/$\textgreater
% \\ Removed by DocsPostProcessor.
 
 Автоматически управляемый хоккеист. Самостоятельно перемещается вдоль вертикального отрезка,
 расположенного перед воротами. Старается держаться на одной горизонтальной линии с шайбой.
 Скорость вратаря ограничена значением {\tt game.goalieMaxSpeed}.
 
% Removed by DocsPostProcessor:
% \textless p$/$\textgreater
% \\ Removed by DocsPostProcessor.
 
 Вратарь единственный среди хоккеистов взаимодействует с шайбой напрямую (как физический объект).}\end{itemize}
}
\item{
public static final HockeyistType VERSATILE\begin{itemize}\item{\vskip -.9ex Хоккеист со сбалансированными параметрами.}\end{itemize}
}
\item{
public static final HockeyistType FORWARD\begin{itemize}\item{\vskip -.9ex Нападающий.
 
% Removed by DocsPostProcessor:
% \textless p$/$\textgreater
% \\ Removed by DocsPostProcessor.
 
 Хоккеист, основным достоинством которого является сила удара.
 Его слабой стороной является стойкость.}\end{itemize}
}
\item{
public static final HockeyistType DEFENCEMAN\begin{itemize}\item{\vskip -.9ex Защитник.
 
% Removed by DocsPostProcessor:
% \textless p$/$\textgreater
% \\ Removed by DocsPostProcessor.
 
 Хоккеист, основным достоинством которого является стойкость.
 Его слабой стороной является точность удара.}\end{itemize}
}
\item{
public static final HockeyistType RANDOM\begin{itemize}\item{\vskip -.9ex Хоккеист со случайными параметрами.}\end{itemize}
}
\end{itemize}
}
\startsubsubsection{Methods}{
\vskip -2em
\begin{itemize}
\item{\vskip -1.9ex 
\membername{valueOf}
{\tt public static HockeyistType {\bf valueOf}( {\tt String } {\bf name} )
\label{l195}\label{l196}}%end signature
}%end item
\divideents{values}
\item{\vskip -1.9ex 
\membername{values}
{\tt public static HockeyistType[] {\bf values}(  )
\label{l197}\label{l198}}%end signature
}%end item
\end{itemize}
}
\hide{inherited}{
\startsubsubsection{Methods inherited from class {\tt Enum}}{
\par{\small 
\refdefined{l14}\vskip -2em
\begin{itemize}
\item{\vskip -1.9ex 
\membername{clone}
{\tt protected final Object {\bf clone}(  )
}%end signature
}%end item
\divideents{compareTo}
\item{\vskip -1.9ex 
\membername{compareTo}
{\tt public final int {\bf compareTo}( {\tt Enum } {\bf arg0} )
}%end signature
}%end item
\divideents{equals}
\item{\vskip -1.9ex 
\membername{equals}
{\tt public final boolean {\bf equals}( {\tt Object } {\bf arg0} )
}%end signature
}%end item
\divideents{finalize}
\item{\vskip -1.9ex 
\membername{finalize}
{\tt protected final void {\bf finalize}(  )
}%end signature
}%end item
\divideents{getDeclaringClass}
\item{\vskip -1.9ex 
\membername{getDeclaringClass}
{\tt public final Class {\bf getDeclaringClass}(  )
}%end signature
}%end item
\divideents{hashCode}
\item{\vskip -1.9ex 
\membername{hashCode}
{\tt public final int {\bf hashCode}(  )
}%end signature
}%end item
\divideents{name}
\item{\vskip -1.9ex 
\membername{name}
{\tt public final String {\bf name}(  )
}%end signature
}%end item
\divideents{ordinal}
\item{\vskip -1.9ex 
\membername{ordinal}
{\tt public final int {\bf ordinal}(  )
}%end signature
}%end item
\divideents{toString}
\item{\vskip -1.9ex 
\membername{toString}
{\tt public String {\bf toString}(  )
}%end signature
}%end item
\divideents{valueOf}
\item{\vskip -1.9ex 
\membername{valueOf}
{\tt public static Enum {\bf valueOf}( {\tt Class } {\bf arg0},
{\tt String } {\bf arg1} )
}%end signature
}%end item
\end{itemize}
}}
}
}
\startsection{Class}{Move}{l5}{%
{\small Стратегия игрока может управлять хоккеистом посредством установки свойств объекта данного класса.}
\vskip .1in 
\startsubsubsection{Declaration}{
\fbox{\vbox{
\hbox{\vbox{\small public 
class 
Move}}
\noindent\hbox{\vbox{{\bf extends} Object}}
}}}

% Removed by DocsPostProcessor:
% \startsubsubsection{Constructors}{
% \vskip -2em
% \begin{itemize}
% \item{\vskip -1.9ex 
% \membername{Move}
% {\tt public {\bf Move}(  )
% \label{l199}\label{l200}}%end signature
% }%end item
% \end{itemize}
% }
% \\ Removed by DocsPostProcessor.

\startsubsubsection{Methods}{
\vskip -2em
\begin{itemize}
\item{\vskip -1.9ex 
\membername{getAction}
{\tt public ActionType {\bf getAction}(  )
\label{l201}\label{l202}}%end signature
\begin{itemize}
\sld
\item{{\bf Returns} - 
Возвращает текущее действие хоккеиста. 
}%end item
\end{itemize}
}%end item
\divideents{getPassAngle}
\item{\vskip -1.9ex 
\membername{getPassAngle}
{\tt public double {\bf getPassAngle}(  )
\label{l203}\label{l204}}%end signature
\begin{itemize}
\sld
\item{{\bf Returns} - 
Возвращает текущее направление паса. 
}%end item
\end{itemize}
}%end item
\divideents{getPassPower}
\item{\vskip -1.9ex 
\membername{getPassPower}
{\tt public double {\bf getPassPower}(  )
\label{l205}\label{l206}}%end signature
\begin{itemize}
\sld
\item{{\bf Returns} - 
Возвращает текущую силу паса. 
}%end item
\end{itemize}
}%end item
\divideents{getSpeedUp}
\item{\vskip -1.9ex 
\membername{getSpeedUp}
{\tt public double {\bf getSpeedUp}(  )
\label{l207}\label{l208}}%end signature
\begin{itemize}
\sld
\item{{\bf Returns} - 
Возвращает текущее ускорение хоккеиста. 
}%end item
\end{itemize}
}%end item
\divideents{getTeammateIndex}
\item{\vskip -1.9ex 
\membername{getTeammateIndex}
{\tt public int {\bf getTeammateIndex}(  )
\label{l209}\label{l210}}%end signature
\begin{itemize}
\sld
\item{{\bf Returns} - 
Возвращает текущий индекс хоккеиста, на которого будет произведена замена,
         или {\tt -1}, если хоккеист не был указан. 
}%end item
\end{itemize}
}%end item
\divideents{getTurn}
\item{\vskip -1.9ex 
\membername{getTurn}
{\tt public double {\bf getTurn}(  )
\label{l211}\label{l212}}%end signature
\begin{itemize}
\sld
\item{{\bf Returns} - 
Возвращает текущий угол поворота хоккеиста. 
}%end item
\end{itemize}
}%end item
\divideents{setAction}
\item{\vskip -1.9ex 
\membername{setAction}
{\tt public void {\bf setAction}( {\tt ActionType } {\bf action} )
\label{l213}\label{l214}}%end signature
\begin{itemize}
\sld
\item{
\sld
{\bf Usage}
  \begin{itemize}\isep
   \item{
Устанавливает действие хоккеиста.
}%end item
  \end{itemize}
}
\end{itemize}
}%end item
\divideents{setPassAngle}
\item{\vskip -1.9ex 
\membername{setPassAngle}
{\tt public void {\bf setPassAngle}( {\tt double } {\bf passAngle} )
\label{l215}\label{l216}}%end signature
\begin{itemize}
\sld
\item{
\sld
{\bf Usage}
  \begin{itemize}\isep
   \item{
Устанавливает направление паса ({\tt ActionType.PASS}).
 
% Removed by DocsPostProcessor:
% \textless p$/$\textgreater
% \\ Removed by DocsPostProcessor.
 
 Направление паса задаётся в радианах относительно текущего направления хоккеиста
 и должно лежать в интервале от {\tt -0.5 * game.passSector} до {\tt 0.5 * game.passSector}.
 Значения, выходящие за указанный интервал, будут приведены к ближайшей его границе.
}%end item
  \end{itemize}
}
\end{itemize}
}%end item
\divideents{setPassPower}
\item{\vskip -1.9ex 
\membername{setPassPower}
{\tt public void {\bf setPassPower}( {\tt double } {\bf passPower} )
\label{l217}\label{l218}}%end signature
\begin{itemize}
\sld
\item{
\sld
{\bf Usage}
  \begin{itemize}\isep
   \item{
Устанавливает силу паса ({\tt ActionType.PASS}).
 
% Removed by DocsPostProcessor:
% \textless p$/$\textgreater
% \\ Removed by DocsPostProcessor.
 
 Сила паса является относительной величиной и должна лежать в интервале от {\tt 0.0} до {\tt 1.0}.
 Значения, выходящие за указанный интервал, будут приведены к ближайшей его границе.
 К значению реальной силы паса применяется также поправочный коэффициент {\tt game.passPowerFactor}.
}%end item
  \end{itemize}
}
\end{itemize}
}%end item
\divideents{setSpeedUp}
\item{\vskip -1.9ex 
\membername{setSpeedUp}
{\tt public void {\bf setSpeedUp}( {\tt double } {\bf speedUp} )
\label{l219}\label{l220}}%end signature
\begin{itemize}
\sld
\item{
\sld
{\bf Usage}
  \begin{itemize}\isep
   \item{
Устанавливает ускорение хоккеиста.
 
% Removed by DocsPostProcessor:
% \textless p$/$\textgreater
% \\ Removed by DocsPostProcessor.
 
 Ускорение является относительным и должно лежать в интервале от {\tt -1.0} до {\tt 1.0}.
 Значения, выходящие за указанный интервал, будут приведены к ближайшей его границе.
}%end item
  \end{itemize}
}
\end{itemize}
}%end item
\divideents{setTeammateIndex}
\item{\vskip -1.9ex 
\membername{setTeammateIndex}
{\tt public void {\bf setTeammateIndex}( {\tt int } {\bf teammateIndex} )
\label{l221}\label{l222}}%end signature
\begin{itemize}
\sld
\item{
\sld
{\bf Usage}
  \begin{itemize}\isep
   \item{
Устанавливает индекс хоккеиста для выполнения замены ({\tt ActionType.SUBSTITUTE}).
 
% Removed by DocsPostProcessor:
% \textless p$/$\textgreater
% \\ Removed by DocsPostProcessor.
 
 Индексация начинается с нуля. Значением по умолчанию является {\tt -1}.
 Если в команде игрока не существует хоккеиста с указанным индексом, то замена произведена не будет.
}%end item
  \end{itemize}
}
\end{itemize}
}%end item
\divideents{setTurn}
\item{\vskip -1.9ex 
\membername{setTurn}
{\tt public void {\bf setTurn}( {\tt double } {\bf turn} )
\label{l223}\label{l224}}%end signature
\begin{itemize}
\sld
\item{
\sld
{\bf Usage}
  \begin{itemize}\isep
   \item{
Устанавливает угол поворота хоккеиста.
 
% Removed by DocsPostProcessor:
% \textless p$/$\textgreater
% \\ Removed by DocsPostProcessor.
 
 Угол поворота задаётся в радианах относительно текущего направления хоккеиста и для хоккеиста
 с базовым значением атрибута подвижность и максимальным запасом выносливости ограничен
 интервалом от {\tt -game.hockeyistTurnAngleFactor} до {\tt game.hockeyistTurnAngleFactor}.
 Значения, выходящие за указанный интервал, будут приведены к ближайшей его границе.
 Положительные значения соответствуют повороту по часовой стрелке.
}%end item
  \end{itemize}
}
\end{itemize}
}%end item
\end{itemize}
}
\hide{inherited}{
}
}
\startsection{Class}{Player}{l6}{%
{\small Содержит данные о текущем состоянии игрока.}
\vskip .1in 
\startsubsubsection{Declaration}{
\fbox{\vbox{
\hbox{\vbox{\small public 
class 
Player}}
\noindent\hbox{\vbox{{\bf extends} Object}}
}}}

% Removed by DocsPostProcessor:
% \startsubsubsection{Constructors}{
% \vskip -2em
% \begin{itemize}
% \item{\vskip -1.9ex 
% \membername{Player}
% {\tt public {\bf Player}( {\tt long } {\bf id},
% {\tt boolean } {\bf me},
% {\tt String } {\bf name},
% {\tt int } {\bf goalCount},
% {\tt boolean } {\bf strategyCrashed},
% {\tt double } {\bf netTop},
% {\tt double } {\bf netLeft},
% {\tt double } {\bf netBottom},
% {\tt double } {\bf netRight},
% {\tt double } {\bf netFront},
% {\tt double } {\bf netBack},
% {\tt boolean } {\bf justScoredGoal},
% {\tt boolean } {\bf justMissedGoal} )
% \label{l225}\label{l226}}%end signature
% }%end item
% \end{itemize}
% }
% \\ Removed by DocsPostProcessor.

\startsubsubsection{Methods}{
\vskip -2em
\begin{itemize}
\item{\vskip -1.9ex 
\membername{getGoalCount}
{\tt public int {\bf getGoalCount}(  )
\label{l227}\label{l228}}%end signature
\begin{itemize}
\sld
\item{{\bf Returns} - 
Возвращает количество шайб, заброшенных хоккеистами данного игрока в сетку противника.
         Шайбы, заброшенные во время состояния вне игры, не влияют на этот счётчик. 
}%end item
\end{itemize}
}%end item
\divideents{getId}
\item{\vskip -1.9ex 
\membername{getId}
{\tt public long {\bf getId}(  )
\label{l229}\label{l230}}%end signature
\begin{itemize}
\sld
\item{{\bf Returns} - 
Возвращает уникальный идентификатор игрока. 
}%end item
\end{itemize}
}%end item
\divideents{getName}
\item{\vskip -1.9ex 
\membername{getName}
{\tt public String {\bf getName}(  )
\label{l231}\label{l232}}%end signature
\begin{itemize}
\sld
\item{{\bf Returns} - 
Возвращает имя игрока. 
}%end item
\end{itemize}
}%end item
\divideents{getNetBack}
\item{\vskip -1.9ex 
\membername{getNetBack}
{\tt public double {\bf getNetBack}(  )
\label{l233}\label{l234}}%end signature
\begin{itemize}
\sld
\item{{\bf Returns} - 
Возвращает абсциссу дальней от вратаря вертикальной границы ворот.
         Соответствует одному из значений {\tt netLeft} или {\tt netRight}. 
}%end item
\end{itemize}
}%end item
\divideents{getNetBottom}
\item{\vskip -1.9ex 
\membername{getNetBottom}
{\tt public double {\bf getNetBottom}(  )
\label{l235}\label{l236}}%end signature
\begin{itemize}
\sld
\item{{\bf Returns} - 
Возвращает ординату нижней штанги ворот. 
}%end item
\end{itemize}
}%end item
\divideents{getNetFront}
\item{\vskip -1.9ex 
\membername{getNetFront}
{\tt public double {\bf getNetFront}(  )
\label{l237}\label{l238}}%end signature
\begin{itemize}
\sld
\item{{\bf Returns} - 
Возвращает абсциссу ближайшей к вратарю вертикальной границы ворот.
         Соответствует одному из значений {\tt netLeft} или {\tt netRight}. 
}%end item
\end{itemize}
}%end item
\divideents{getNetLeft}
\item{\vskip -1.9ex 
\membername{getNetLeft}
{\tt public double {\bf getNetLeft}(  )
\label{l239}\label{l240}}%end signature
\begin{itemize}
\sld
\item{{\bf Returns} - 
Возвращает абсциссу левой границы ворот. 
}%end item
\end{itemize}
}%end item
\divideents{getNetRight}
\item{\vskip -1.9ex 
\membername{getNetRight}
{\tt public double {\bf getNetRight}(  )
\label{l241}\label{l242}}%end signature
\begin{itemize}
\sld
\item{{\bf Returns} - 
Возвращает абсциссу правой границы ворот. 
}%end item
\end{itemize}
}%end item
\divideents{getNetTop}
\item{\vskip -1.9ex 
\membername{getNetTop}
{\tt public double {\bf getNetTop}(  )
\label{l243}\label{l244}}%end signature
\begin{itemize}
\sld
\item{{\bf Returns} - 
Возвращает ординату верхней штанги ворот. 
}%end item
\end{itemize}
}%end item
\divideents{isJustMissedGoal}
\item{\vskip -1.9ex 
\membername{isJustMissedGoal}
{\tt public boolean {\bf isJustMissedGoal}(  )
\label{l245}\label{l246}}%end signature
\begin{itemize}
\sld
\item{{\bf Returns} - 
Возвращает {\tt true} в том и только в том случае, если игрок только что пропустил гол.
         
% Removed by DocsPostProcessor:
% \textless p$/$\textgreater
% \\ Removed by DocsPostProcessor.
 
         Вместе с установленным флагом {\tt justScoredGoal} другого игрока означает,
         что сейчас состояние вне игры и новые голы не будут засчитаны.
         Длительность состояния вне игры составляет {\tt game.afterGoalStateTickCount} тиков. 
}%end item
\end{itemize}
}%end item
\divideents{isJustScoredGoal}
\item{\vskip -1.9ex 
\membername{isJustScoredGoal}
{\tt public boolean {\bf isJustScoredGoal}(  )
\label{l247}\label{l248}}%end signature
\begin{itemize}
\sld
\item{{\bf Returns} - 
Возвращает {\tt true} в том и только в том случае, если игрок только что забил гол.
         
% Removed by DocsPostProcessor:
% \textless p$/$\textgreater
% \\ Removed by DocsPostProcessor.
 
         Вместе с установленным флагом {\tt justMissedGoal} другого игрока означает,
         что сейчас состояние вне игры и новые голы не будут засчитаны.
         Длительность состояния вне игры составляет {\tt game.afterGoalStateTickCount} тиков. 
}%end item
\end{itemize}
}%end item
\divideents{isMe}
\item{\vskip -1.9ex 
\membername{isMe}
{\tt public boolean {\bf isMe}(  )
\label{l249}\label{l250}}%end signature
\begin{itemize}
\sld
\item{{\bf Returns} - 
Возвращает {\tt true} в том и только в том случае, если этот игрок ваш. 
}%end item
\end{itemize}
}%end item
\divideents{isStrategyCrashed}
\item{\vskip -1.9ex 
\membername{isStrategyCrashed}
{\tt public boolean {\bf isStrategyCrashed}(  )
\label{l251}\label{l252}}%end signature
\begin{itemize}
\sld
\item{{\bf Returns} - 
Возвращает специальный флаг --- показатель того, что стратегия игрока \textless \textless упала\textgreater \textgreater .
         Более подробную информацию можно найти в документации к игре. 
}%end item
\end{itemize}
}%end item
\end{itemize}
}
\hide{inherited}{
}
}
\startsection{Class}{Puck}{l7}{%
{\small Класс, определяющий хоккейную шайбу. Содержит также все свойства юнита.}
\vskip .1in 
\startsubsubsection{Declaration}{
\fbox{\vbox{
\hbox{\vbox{\small public 
class 
Puck}}
\noindent\hbox{\vbox{{\bf extends} Unit}}
}}}

% Removed by DocsPostProcessor:
% \startsubsubsection{Constructors}{
% \vskip -2em
% \begin{itemize}
% \item{\vskip -1.9ex 
% \membername{Puck}
% {\tt public {\bf Puck}( {\tt long } {\bf id},
% {\tt double } {\bf mass},
% {\tt double } {\bf radius},
% {\tt double } {\bf x},
% {\tt double } {\bf y},
% {\tt double } {\bf speedX},
% {\tt double } {\bf speedY},
% {\tt long } {\bf ownerHockeyistId},
% {\tt long } {\bf ownerPlayerId} )
% \label{l253}\label{l254}}%end signature
% }%end item
% \end{itemize}
% }
% \\ Removed by DocsPostProcessor.

\startsubsubsection{Methods}{
\vskip -2em
\begin{itemize}
\item{\vskip -1.9ex 
\membername{getOwnerHockeyistId}
{\tt public long {\bf getOwnerHockeyistId}(  )
\label{l255}\label{l256}}%end signature
\begin{itemize}
\sld
\item{{\bf Returns} - 
Возвращает идентификатор хоккеиста, контролирующего шайбу, или {\tt -1}. 
}%end item
\end{itemize}
}%end item
\divideents{getOwnerPlayerId}
\item{\vskip -1.9ex 
\membername{getOwnerPlayerId}
{\tt public long {\bf getOwnerPlayerId}(  )
\label{l257}\label{l258}}%end signature
\begin{itemize}
\sld
\item{{\bf Returns} - 
Возвращает идентификатор игрока, контролирующего шайбу, или {\tt -1}. 
}%end item
\end{itemize}
}%end item
\end{itemize}
}
\hide{inherited}{
\startsubsubsection{Methods inherited from class {\tt Unit}}{
\par{\small 
\refdefined{l8}\vskip -2em
\begin{itemize}
\item{\vskip -1.9ex 
\membername{getAngle}
{\tt public final double {\bf getAngle}(  )
}%end signature
\begin{itemize}
\sld
\item{{\bf Returns} - 
Возвращает угол поворота объекта в радианах. Нулевой угол соответствует направлению оси абсцисс.
         Положительные значения соответствуют повороту по часовой стрелке. 
}%end item
\end{itemize}
}%end item
\divideents{getAngleTo}
\item{\vskip -1.9ex 
\membername{getAngleTo}
{\tt public double {\bf getAngleTo}( {\tt double } {\bf x},
{\tt double } {\bf y} )
}%end signature
\begin{itemize}
\sld
\item{
\sld
{\bf Parameters}
\sld\isep
  \begin{itemize}
\sld\isep
   \item{
\sld
{\tt x} - X-координата точки.}
   \item{
\sld
{\tt y} - Y-координата точки.}
  \end{itemize}
}%end item
\item{{\bf Returns} - 
Возвращает ориентированный угол [{\tt -PI}, {\tt PI}] между направлением
         данного объекта и вектором из центра данного объекта к указанной точке. 
}%end item
\end{itemize}
}%end item
\divideents{getAngleTo}
\item{\vskip -1.9ex 
\membername{getAngleTo}
{\tt public double {\bf getAngleTo}( {\tt Unit } {\bf unit} )
}%end signature
\begin{itemize}
\sld
\item{
\sld
{\bf Parameters}
\sld\isep
  \begin{itemize}
\sld\isep
   \item{
\sld
{\tt unit} - Объект, к центру которого необходимо определить угол.}
  \end{itemize}
}%end item
\item{{\bf Returns} - 
Возвращает ориентированный угол [{\tt -PI}, {\tt PI}] между направлением
         данного объекта и вектором из центра данного объекта к центру указанного объекта. 
}%end item
\end{itemize}
}%end item
\divideents{getAngularSpeed}
\item{\vskip -1.9ex 
\membername{getAngularSpeed}
{\tt public double {\bf getAngularSpeed}(  )
}%end signature
\begin{itemize}
\sld
\item{{\bf Returns} - 
Возвращает скорость вращения объекта.
         Положительные значения соответствуют вращению по часовой стрелке. 
}%end item
\end{itemize}
}%end item
\divideents{getDistanceTo}
\item{\vskip -1.9ex 
\membername{getDistanceTo}
{\tt public double {\bf getDistanceTo}( {\tt double } {\bf x},
{\tt double } {\bf y} )
}%end signature
\begin{itemize}
\sld
\item{
\sld
{\bf Parameters}
\sld\isep
  \begin{itemize}
\sld\isep
   \item{
\sld
{\tt x} - X-координата точки.}
   \item{
\sld
{\tt y} - Y-координата точки.}
  \end{itemize}
}%end item
\item{{\bf Returns} - 
Возвращает расстояние до точки от центра данного объекта. 
}%end item
\end{itemize}
}%end item
\divideents{getDistanceTo}
\item{\vskip -1.9ex 
\membername{getDistanceTo}
{\tt public double {\bf getDistanceTo}( {\tt Unit } {\bf unit} )
}%end signature
\begin{itemize}
\sld
\item{
\sld
{\bf Parameters}
\sld\isep
  \begin{itemize}
\sld\isep
   \item{
\sld
{\tt unit} - Объект, до центра которого необходимо определить расстояние.}
  \end{itemize}
}%end item
\item{{\bf Returns} - 
Возвращает расстояние от центра данного объекта до центра указанного объекта. 
}%end item
\end{itemize}
}%end item
\divideents{getId}
\item{\vskip -1.9ex 
\membername{getId}
{\tt public long {\bf getId}(  )
}%end signature
\begin{itemize}
\sld
\item{{\bf Returns} - 
Возвращает уникальный идентификатор объекта. 
}%end item
\end{itemize}
}%end item
\divideents{getMass}
\item{\vskip -1.9ex 
\membername{getMass}
{\tt public double {\bf getMass}(  )
}%end signature
\begin{itemize}
\sld
\item{{\bf Returns} - 
Возвращает массу объекта в единицах массы. 
}%end item
\end{itemize}
}%end item
\divideents{getRadius}
\item{\vskip -1.9ex 
\membername{getRadius}
{\tt public double {\bf getRadius}(  )
}%end signature
\begin{itemize}
\sld
\item{{\bf Returns} - 
Возвращает радиус объекта. 
}%end item
\end{itemize}
}%end item
\divideents{getSpeedX}
\item{\vskip -1.9ex 
\membername{getSpeedX}
{\tt public final double {\bf getSpeedX}(  )
}%end signature
\begin{itemize}
\sld
\item{{\bf Returns} - 
Возвращает X-составляющую скорости объекта. Ось абсцисс направлена слева направо. 
}%end item
\end{itemize}
}%end item
\divideents{getSpeedY}
\item{\vskip -1.9ex 
\membername{getSpeedY}
{\tt public final double {\bf getSpeedY}(  )
}%end signature
\begin{itemize}
\sld
\item{{\bf Returns} - 
Возвращает Y-составляющую скорости объекта. Ось ординат направлена свеху вниз. 
}%end item
\end{itemize}
}%end item
\divideents{getX}
\item{\vskip -1.9ex 
\membername{getX}
{\tt public final double {\bf getX}(  )
}%end signature
\begin{itemize}
\sld
\item{{\bf Returns} - 
Возвращает X-координату центра объекта. Ось абсцисс направлена слева направо. 
}%end item
\end{itemize}
}%end item
\divideents{getY}
\item{\vskip -1.9ex 
\membername{getY}
{\tt public final double {\bf getY}(  )
}%end signature
\begin{itemize}
\sld
\item{{\bf Returns} - 
Возвращает Y-координату центра объекта. Ось ординат направлена свеху вниз. 
}%end item
\end{itemize}
}%end item
\end{itemize}
}}
}
}
\startsection{Class}{Unit}{l8}{%
{\small Базовый класс для определения объектов (\textless \textless юнитов\textgreater \textgreater ) на игровом поле.}
\vskip .1in 
\startsubsubsection{Declaration}{
\fbox{\vbox{
\hbox{\vbox{\small public abstract 
class 
Unit}}
\noindent\hbox{\vbox{{\bf extends} Object}}
}}}

% Removed by DocsPostProcessor:
% \startsubsubsection{Constructors}{
% \vskip -2em
% \begin{itemize}
% \item{\vskip -1.9ex 
% \membername{Unit}
% {\tt protected {\bf Unit}( {\tt long } {\bf id},
% {\tt double } {\bf mass},
% {\tt double } {\bf radius},
% {\tt double } {\bf x},
% {\tt double } {\bf y},
% {\tt double } {\bf speedX},
% {\tt double } {\bf speedY},
% {\tt double } {\bf angle},
% {\tt double } {\bf angularSpeed} )
% \label{l259}\label{l260}}%end signature
% }%end item
% \end{itemize}
% }
% \\ Removed by DocsPostProcessor.

\startsubsubsection{Methods}{
\vskip -2em
\begin{itemize}
\item{\vskip -1.9ex 
\membername{getAngle}
{\tt public final double {\bf getAngle}(  )
\label{l261}\label{l262}}%end signature
\begin{itemize}
\sld
\item{{\bf Returns} - 
Возвращает угол поворота объекта в радианах. Нулевой угол соответствует направлению оси абсцисс.
         Положительные значения соответствуют повороту по часовой стрелке. 
}%end item
\end{itemize}
}%end item
\divideents{getAngleTo}
\item{\vskip -1.9ex 
\membername{getAngleTo}
{\tt public double {\bf getAngleTo}( {\tt double } {\bf x},
{\tt double } {\bf y} )
\label{l263}\label{l264}}%end signature
\begin{itemize}
\sld
\item{
\sld
{\bf Parameters}
\sld\isep
  \begin{itemize}
\sld\isep
   \item{
\sld
{\tt x} - X-координата точки.}
   \item{
\sld
{\tt y} - Y-координата точки.}
  \end{itemize}
}%end item
\item{{\bf Returns} - 
Возвращает ориентированный угол [{\tt -PI}, {\tt PI}] между направлением
         данного объекта и вектором из центра данного объекта к указанной точке. 
}%end item
\end{itemize}
}%end item
\divideents{getAngleTo}
\item{\vskip -1.9ex 
\membername{getAngleTo}
{\tt public double {\bf getAngleTo}( {\tt Unit } {\bf unit} )
\label{l265}\label{l266}}%end signature
\begin{itemize}
\sld
\item{
\sld
{\bf Parameters}
\sld\isep
  \begin{itemize}
\sld\isep
   \item{
\sld
{\tt unit} - Объект, к центру которого необходимо определить угол.}
  \end{itemize}
}%end item
\item{{\bf Returns} - 
Возвращает ориентированный угол [{\tt -PI}, {\tt PI}] между направлением
         данного объекта и вектором из центра данного объекта к центру указанного объекта. 
}%end item
\end{itemize}
}%end item
\divideents{getAngularSpeed}
\item{\vskip -1.9ex 
\membername{getAngularSpeed}
{\tt public double {\bf getAngularSpeed}(  )
\label{l267}\label{l268}}%end signature
\begin{itemize}
\sld
\item{{\bf Returns} - 
Возвращает скорость вращения объекта.
         Положительные значения соответствуют вращению по часовой стрелке. 
}%end item
\end{itemize}
}%end item
\divideents{getDistanceTo}
\item{\vskip -1.9ex 
\membername{getDistanceTo}
{\tt public double {\bf getDistanceTo}( {\tt double } {\bf x},
{\tt double } {\bf y} )
\label{l269}\label{l270}}%end signature
\begin{itemize}
\sld
\item{
\sld
{\bf Parameters}
\sld\isep
  \begin{itemize}
\sld\isep
   \item{
\sld
{\tt x} - X-координата точки.}
   \item{
\sld
{\tt y} - Y-координата точки.}
  \end{itemize}
}%end item
\item{{\bf Returns} - 
Возвращает расстояние до точки от центра данного объекта. 
}%end item
\end{itemize}
}%end item
\divideents{getDistanceTo}
\item{\vskip -1.9ex 
\membername{getDistanceTo}
{\tt public double {\bf getDistanceTo}( {\tt Unit } {\bf unit} )
\label{l271}\label{l272}}%end signature
\begin{itemize}
\sld
\item{
\sld
{\bf Parameters}
\sld\isep
  \begin{itemize}
\sld\isep
   \item{
\sld
{\tt unit} - Объект, до центра которого необходимо определить расстояние.}
  \end{itemize}
}%end item
\item{{\bf Returns} - 
Возвращает расстояние от центра данного объекта до центра указанного объекта. 
}%end item
\end{itemize}
}%end item
\divideents{getId}
\item{\vskip -1.9ex 
\membername{getId}
{\tt public long {\bf getId}(  )
\label{l273}\label{l274}}%end signature
\begin{itemize}
\sld
\item{{\bf Returns} - 
Возвращает уникальный идентификатор объекта. 
}%end item
\end{itemize}
}%end item
\divideents{getMass}
\item{\vskip -1.9ex 
\membername{getMass}
{\tt public double {\bf getMass}(  )
\label{l275}\label{l276}}%end signature
\begin{itemize}
\sld
\item{{\bf Returns} - 
Возвращает массу объекта в единицах массы. 
}%end item
\end{itemize}
}%end item
\divideents{getRadius}
\item{\vskip -1.9ex 
\membername{getRadius}
{\tt public double {\bf getRadius}(  )
\label{l277}\label{l278}}%end signature
\begin{itemize}
\sld
\item{{\bf Returns} - 
Возвращает радиус объекта. 
}%end item
\end{itemize}
}%end item
\divideents{getSpeedX}
\item{\vskip -1.9ex 
\membername{getSpeedX}
{\tt public final double {\bf getSpeedX}(  )
\label{l279}\label{l280}}%end signature
\begin{itemize}
\sld
\item{{\bf Returns} - 
Возвращает X-составляющую скорости объекта. Ось абсцисс направлена слева направо. 
}%end item
\end{itemize}
}%end item
\divideents{getSpeedY}
\item{\vskip -1.9ex 
\membername{getSpeedY}
{\tt public final double {\bf getSpeedY}(  )
\label{l281}\label{l282}}%end signature
\begin{itemize}
\sld
\item{{\bf Returns} - 
Возвращает Y-составляющую скорости объекта. Ось ординат направлена свеху вниз. 
}%end item
\end{itemize}
}%end item
\divideents{getX}
\item{\vskip -1.9ex 
\membername{getX}
{\tt public final double {\bf getX}(  )
\label{l283}\label{l284}}%end signature
\begin{itemize}
\sld
\item{{\bf Returns} - 
Возвращает X-координату центра объекта. Ось абсцисс направлена слева направо. 
}%end item
\end{itemize}
}%end item
\divideents{getY}
\item{\vskip -1.9ex 
\membername{getY}
{\tt public final double {\bf getY}(  )
\label{l285}\label{l286}}%end signature
\begin{itemize}
\sld
\item{{\bf Returns} - 
Возвращает Y-координату центра объекта. Ось ординат направлена свеху вниз. 
}%end item
\end{itemize}
}%end item
\end{itemize}
}
\hide{inherited}{
}
}
\startsection{Class}{World}{l9}{%
{\small Этот класс описывает игровой мир. Содержит также описания всех игроков и игровых объектов (\textless \textless юнитов\textgreater \textgreater ).}
\vskip .1in 
\startsubsubsection{Declaration}{
\fbox{\vbox{
\hbox{\vbox{\small public 
class 
World}}
\noindent\hbox{\vbox{{\bf extends} Object}}
}}}

% Removed by DocsPostProcessor:
% \startsubsubsection{Constructors}{
% \vskip -2em
% \begin{itemize}
% \item{\vskip -1.9ex 
% \membername{World}
% {\tt public {\bf World}( {\tt int } {\bf tick},
% {\tt int } {\bf tickCount},
% {\tt double } {\bf width},
% {\tt double } {\bf height},
% {\tt Player[]} {\bf players},
% {\tt Hockeyist[]} {\bf hockeyists},
% {\tt Puck } {\bf puck} )
% \label{l287}\label{l288}}%end signature
% }%end item
% \end{itemize}
% }
% \\ Removed by DocsPostProcessor.

\startsubsubsection{Methods}{
\vskip -2em
\begin{itemize}
\item{\vskip -1.9ex 
\membername{getHeight}
{\tt public double {\bf getHeight}(  )
\label{l289}\label{l290}}%end signature
\begin{itemize}
\sld
\item{{\bf Returns} - 
Возвращает высоту мира. 
}%end item
\end{itemize}
}%end item
\divideents{getHockeyists}
\item{\vskip -1.9ex 
\membername{getHockeyists}
{\tt public Hockeyist[] {\bf getHockeyists}(  )
\label{l291}\label{l292}}%end signature
\begin{itemize}
\sld
\item{{\bf Returns} - 
Возвращает список хоккеистов (в случайном порядке), включая вратарей и хоккеиста стратегии,
         вызвавшей этот метод. После каждого тика объекты, задающие хоккеистов, пересоздаются. 
}%end item
\end{itemize}
}%end item
\divideents{getMyPlayer}
\item{\vskip -1.9ex 
\membername{getMyPlayer}
{\tt public Player {\bf getMyPlayer}(  )
\label{l293}\label{l294}}%end signature
\begin{itemize}
\sld
\item{{\bf Returns} - 
Возвращает вашего игрока. 
}%end item
\end{itemize}
}%end item
\divideents{getOpponentPlayer}
\item{\vskip -1.9ex 
\membername{getOpponentPlayer}
{\tt public Player {\bf getOpponentPlayer}(  )
\label{l295}\label{l296}}%end signature
\begin{itemize}
\sld
\item{{\bf Returns} - 
Возвращает игрока, соревнующегося с вами. 
}%end item
\end{itemize}
}%end item
\divideents{getPlayers}
\item{\vskip -1.9ex 
\membername{getPlayers}
{\tt public Player[] {\bf getPlayers}(  )
\label{l297}\label{l298}}%end signature
\begin{itemize}
\sld
\item{{\bf Returns} - 
Возвращает список игроков (в случайном порядке).
         После каждого тика объекты, задающие игроков, пересоздаются. 
}%end item
\end{itemize}
}%end item
\divideents{getPuck}
\item{\vskip -1.9ex 
\membername{getPuck}
{\tt public Puck {\bf getPuck}(  )
\label{l299}\label{l300}}%end signature
\begin{itemize}
\sld
\item{{\bf Returns} - 
Возвращает шайбу. 
}%end item
\end{itemize}
}%end item
\divideents{getTick}
\item{\vskip -1.9ex 
\membername{getTick}
{\tt public int {\bf getTick}(  )
\label{l301}\label{l302}}%end signature
\begin{itemize}
\sld
\item{{\bf Returns} - 
Возвращает номер текущего тика. 
}%end item
\end{itemize}
}%end item
\divideents{getTickCount}
\item{\vskip -1.9ex 
\membername{getTickCount}
{\tt public int {\bf getTickCount}(  )
\label{l303}\label{l304}}%end signature
\begin{itemize}
\sld
\item{{\bf Returns} - 
Возвращает базовую длительность игры в тиках.
         Реальная длительность может отличаться от этого значения в большую сторону. 
}%end item
\end{itemize}
}%end item
\divideents{getWidth}
\item{\vskip -1.9ex 
\membername{getWidth}
{\tt public double {\bf getWidth}(  )
\label{l305}\label{l306}}%end signature
\begin{itemize}
\sld
\item{{\bf Returns} - 
Возвращает ширину мира. 
}%end item
\end{itemize}
}%end item
\end{itemize}
}
\hide{inherited}{
}
}
}
}
\newpage
\def\packagename{\textless none\textgreater }
\chapter{\bf Package \textless none\textgreater }{
\vskip -.25in
\hbox to \hsize{\it Package Contents\hfil Page}
\rule{\hsize}{.7mm}
\vskip .13in
\hbox{\bf Interfaces}
\entityintro{Strategy}{l307}{Стратегия --- интерфейс, содержащий описание методов искусственного интеллекта хоккеиста.}
\vskip .1in
\rule{\hsize}{.7mm}
\vskip .1in
\newpage
\section{Interfaces}{
\startsection{Interface}{Strategy}{l307}{%
{\small Стратегия --- интерфейс, содержащий описание методов искусственного интеллекта хоккеиста.
 Каждая пользовательская стратегия должна реализовывать этот интерфейс.
 Может отсутствовать в некоторых языковых пакетах, если язык не поддерживает интерфейсы.}
\vskip .1in 
\startsubsubsection{Declaration}{
\fbox{\vbox{
\hbox{\vbox{\small public interface 
Strategy}}
}}}
\startsubsubsection{Methods}{
\vskip -2em
\begin{itemize}
\item{\vskip -1.9ex 
\membername{move}
{\tt public void {\bf move}( {\tt Hockeyist } {\bf self},
{\tt World } {\bf world},
{\tt Game } {\bf game},
{\tt Move } {\bf move} )
\label{l308}\label{l309}}%end signature
\begin{itemize}
\sld
\item{
\sld
{\bf Usage}
  \begin{itemize}\isep
   \item{
Основной метод стратегии, осуществляющий управление хоккеистом.
 Вызывается каждый тик для каждого хоккеиста.
}%end item
  \end{itemize}
}
\item{
\sld
{\bf Parameters}
\sld\isep
  \begin{itemize}
\sld\isep
   \item{
\sld
{\tt self} - Хоккеист, которым данный метод будет осуществлять управление.}
   \item{
\sld
{\tt world} - Текущее состояние мира.}
   \item{
\sld
{\tt game} - Различные игровые константы.}
   \item{
\sld
{\tt move} - Результатом работы метода является изменение полей данного объекта.}
  \end{itemize}
}%end item
\end{itemize}
}%end item
\end{itemize}
}
\hide{inherited}{
}
}
}
}
\end{document}
